% Préambule commun aux cours IPR Scholar.
\documentclass[11pt,a4paper,oneside]{article}
\usepackage[T1]{fontenc}
\usepackage[utf8]{inputenc}
\usepackage[french]{babel}
\usepackage{lmodern}
\usepackage{microtype}
\usepackage[a4paper,margin=2.35cm]{geometry}
\usepackage{amsmath,amssymb,mathtools}
\usepackage{booktabs,longtable,array}
\usepackage{xcolor}
\usepackage{hyperref}
\usepackage{tikz}
\usetikzlibrary{arrows.meta,positioning,shapes.geometric,calc}
\hypersetup{colorlinks=true,linkcolor=blue!55!black,citecolor=blue!55!black,urlcolor=blue!55!black}
\definecolor{iprblue}{HTML}{1F4E79}
\definecolor{iprgreen}{HTML}{2D6A4F}
\newcommand{\IPR}{\textsc{ipr}}
\newcommand{\coursmeta}[2]{\author{Projet \IPR{} -- documentation Scholar}\date{#1}\maketitle\begin{center}\small #2\end{center}\vspace{0.5cm}\tableofcontents\clearpage}
\newcommand{\resultat}[1]{\paragraph{À retenir.} #1}

\setcounter{tocdepth}{2}
\setlength{\emergencystretch}{2em}
\hypersetup{pdftitle={Cours 11 -- Architecture des images numériques},
pdfauthor={Projet IPR -- Scholar}}
\title{Cours 11 -- Architecture des images numériques\\
\large De la lumière aux octets, des octets à l'image affichée}
\begin{document}
\coursmeta{8 octobre 2026}{Cours de référence : physique, mathématiques,
représentations mémoire, formats, compression, métadonnées et développement RAW.\\
Exemples originaux, travaux pratiques reproductibles et bibliographie commentée.}

\section*{Comment utiliser ce cours}
\addcontentsline{toc}{section}{Comment utiliser ce cours}
Ce cours répond à une question apparemment simple : \emph{que contient un
fichier image et que fait le logiciel pour le montrer ?} La réponse exige
de traverser plusieurs couches. Une mesure de capteur n'est pas une couleur
affichable ; une couleur n'est pas un octet ; un octet n'est pas un fichier ;
un fichier décodable n'est pas nécessairement interprété correctement.

Le parcours débutant suit les parties I, II et V avant les détails des
formats. Le parcours développeur suit les parties II à IV, puis VII.
Pour comprendre les différences Canon/Lightroom/IPR, lire surtout les
parties V et VI, sans omettre la quantification et la gestion des profils.
Les encadrés « À retenir » résument les distinctions essentielles.
Les équations sont accompagnées d'une interprétation et de valeurs numériques.

\paragraph{Prérequis.} Savoir additionner, multiplier, lire un tableau et
comprendre une fonction. Les matrices, logarithmes, probabilités et
transformées sont introduits au moment où ils deviennent utiles. Il ne
s'agit pas d'un catalogue exhaustif de tous les codecs, ni d'une
spécification de remplacement permettant d'écrire seul un décodeur industriel.

\paragraph{Nature des sources.} Les références normatives décrivent des
contrats de fichiers. Les manuels constructeurs décrivent des modèles
et modes précis. Le code des décodeurs et les travaux de rétro-ingénierie
documentent des implémentations, pas une promesse universelle du constructeur.
Les exemples, calculs, diagrammes et exercices sont des constructions
pédagogiques originales. Les liens de la bibliographie ont été recherchés
le 8 octobre 2026 ; certains documents historiques peuvent changer d'adresse.

\paragraph{Périmètre IPR.} Ce cours est un document Scholar, pas une
modification du moteur. Les études Canon locales servent d'illustrations
datées, non de constantes à recopier pour Sony, Nikon ou tous les Canon.
Les espaces de travail, formats et limites d'un logiciel doivent être
vérifiés dans son code et sa matrice de validation.

\part{De la scène à une grille de nombres}

\section{Une image n'est pas seulement un tableau RGB}
\subsection{Quatre niveaux de description}
Une photographie existe successivement comme phénomène lumineux,
mesure électronique, représentation numérique et rendu perceptuel.
Le mot « image » désigne souvent l'un de ces niveaux sans le préciser.

\begin{longtable}{p{3cm}p{5.1cm}p{6.5cm}}
\toprule Niveau & Question & Exemple \\ \midrule
Scène physique & Quelle lumière atteint le capteur ? & Un spectre lumineux,
une durée d'exposition et une distribution spatiale. \\
Signal mesuré & Que mesure chaque photosite ? & Une charge électrique,
puis un code numérique filtré par une mosaïque colorée. \\
Image interprétée & Que signifient les composantes ? & RGB linéaire de
caméra, XYZ, sRGB encodé, profondeur, alpha. \\
Fichier & Comment les valeurs sont-elles stockées ? & TIFF avec offsets,
JPEG avec coefficients, PNG avec blocs compressés. \\
Affichage & Quelle lumière l'écran émet-il ? & Conversion vers le profil
écran, fonction de transfert et limites de luminance. \\ \bottomrule
\end{longtable}

Le même tableau $(128,64,32)$ peut désigner une couleur sRGB, des codes
RGB d'un autre espace ou trois mesures sans signification colorimétrique.
Le même fichier peut contenir plusieurs tableaux : miniature, aperçu,
mosaïque RAW, masque, carte de profondeur ou image HDR auxiliaire.

\resultat{L'architecture complète est : des valeurs, une géométrie,
une signification physique ou colorimétrique, une organisation binaire
et des règles d'interprétation. Il manque de l'information si l'un de
ces éléments est absent.}

\subsection{Image matricielle et image vectorielle}
Une image matricielle est définie sur une grille. Une image vectorielle
décrit des objets : chemins, courbes, texte, surfaces. Un SVG peut dire
« tracer un cercle » ; un PNG doit fournir ou reconstruire les valeurs
de ses pixels. La rasterisation transforme le premier en grille à une
résolution choisie. Un PDF peut mêler texte, vecteurs et photographies.
L'agrandissement d'un vecteur recalcule sa couverture ; l'agrandissement
d'une photo interpole des mesures existantes. Une IA peut inventer une
estimation plausible de détails, pas retrouver systématiquement ceux
qui n'ont jamais été mesurés.

\section{Physique de la formation de l'image}
\subsection{Photons, optique et temps d'exposition}
La lumière peut être décrite par un champ électromagnétique ou par des
photons. L'énergie d'un photon de longueur d'onde $\lambda$ est
\begin{equation}
E_\gamma=\frac{hc}{\lambda},
\end{equation}
où $h$ est la constante de Planck et $c$ la vitesse de la lumière.
À environ $550$ nm, cette énergie vaut $3{,}61\times10^{-19}$ joule.
Un pixel ne stocke pourtant pas une liste de photons : le capteur
intègre leur effet pendant l'exposition.

Un modèle simplifié du nombre moyen d'électrons du photosite $i$ est
\begin{equation}
\mu_i=t\int_{A_i}\int E(x,y,\lambda)\,
\eta_i(\lambda)\,\frac{\lambda}{hc}\,d\lambda\,dA.
\end{equation}
$E$ est ici une irradiance spectrale sur le capteur, $A_i$ l'aire du
photosite, $t$ la durée et $\eta_i$ l'efficacité de conversion. Les filtres
colorés et la transmission optique peuvent être inclus dans $E$ ou
dans une réponse effective $\eta_i$. Ce modèle explique trois faits :
une exposition plus longue augmente le signal, les filtres changent
la réponse spectrale et les photosites ne mesurent pas tous la même
bande de lumière.

Pour une surface diffuse, un modèle de scène souvent utile est
$E(\lambda)\propto S(\lambda)\rho(\lambda)$ : spectre de l'éclairage
fois réflectance. Il n'est pas complet pour les surfaces spéculaires,
translucides ou fluorescentes. Une couleur de peinture publiée sur un
site marchand ne donne ni son spectre de réflectance ni l'éclairage réel.

\subsection{Flou optique et fonction d'étalement}
Même avant la conversion numérique, un point de la scène devient une
tache. Dans un modèle local linéaire invariant, l'image idéale $f$ est
convoluée par une fonction $h$, dite PSF :
\begin{equation}
g(x,y)=(h*f)(x,y)=\iint h(u,v)f(x-u,y-v)\,du\,dv.
\end{equation}
Défocalisation, diffraction et mouvement participent au flou. Dans
le domaine fréquentiel, la convolution devient un produit ; la MTF
décrit la transmission des contrastes suivant la fréquence. Si une
fréquence est supprimée et noyée dans le bruit, une accentuation ne
la recrée pas avec certitude. Elle augmente surtout le contraste de
ce qui subsiste.

\subsection{Exposition et ISO : deux actions différentes}
À scène et géométrie constantes, doubler $t$ double approximativement
les photons collectés tant qu'il n'y a pas saturation. Une variation
de $e$ EV représente un facteur $2^e$. Pour un objectif idéal et un
champ donné, l'éclairement varie approximativement comme $1/N^2$, où
$N$ est le nombre d'ouverture. Passer de $f/4$ à $f/2{,}8$ représente
environ un doublement de l'éclairement.

L'ISO règle une chaîne de gain et de présentation de la sensibilité,
selon l'appareil ; ce n'est pas une commande qui multiplie les photons.
Le gain peut être analogique, numérique ou associé à des modes de
lecture différents. « ISO plus élevé » ne signifie donc pas
automatiquement « moins de photons » : ce sont les choix d'exposition
associés qui le déterminent. Pour comparer le bruit, documenter le
temps, l'ouverture, l'éclairage, l'ISO et le mode RAW.

\section{Capteur, bruit et conversion analogique-numérique}
\subsection{Modèle simplifié d'un code RAW}
Un modèle pédagogique est
\begin{equation}
D_i=\operatorname{clip}\!\left(
\operatorname{round}(B_i+K_i(N_i+n_{r,i})),0,2^b-1\right),
\end{equation}
avec $B_i$ un offset noir en codes, $K_i$ un gain en codes par électron,
$N_i$ les électrons photo-générés, $n_{r,i}$ un bruit de lecture et $b$
la profondeur de conversion. Des traitements internes et des
non-linéarités peuvent compléter ce modèle. La norme de mesure
EMVA~1288 distingue notamment caractérisation linéaire et générale
des capteurs~\cite{emva}. Ce cours utilise le modèle simple pour
raisonner ; il ne certifie aucun boîtier.

\subsection{Bruit de photons et bruit de lecture}
Pour un comptage de photons de type Poisson,
\begin{equation}
N\sim\mathcal P(\mu),\qquad \mathbb E[N]=\mu,
\qquad \operatorname{Var}(N)=\mu.
\end{equation}
L'écart-type vaut $\sqrt\mu$ : recevoir davantage de lumière augmente
le bruit absolu, mais améliore le rapport signal/bruit relatif.
Si le bruit de lecture indépendant a une variance $\sigma_r^2$,
\begin{equation}
\operatorname{SNR}\simeq\frac{\mu}{\sqrt{\mu+\sigma_r^2}}.
\end{equation}
Avec $\mu=100$ électrons et $\sigma_r=3$, on obtient environ $9{,}58$.
Avec $\mu=400$, on obtient $19{,}78$. Multiplier ensuite le fichier
sombre par quatre amplifie aussi son bruit et ne procure pas les
400 électrons d'une vraie exposition plus longue.

Le bruit thermique dépend du temps et de la température ; les
inhomogénéités de gain ou d'offset produisent du bruit spatial fixe.
Un pixel chaud, une ligne de lecture et un bruit aléatoire ne se
traitent pas forcément avec le même algorithme.

\subsection{Niveau noir, niveau blanc et saturation}
Le code zéro n'est pas nécessairement le noir du capteur. Un offset
permet notamment de représenter les fluctuations autour du noir
sans les couper trop tôt. Pour une normalisation simple :
\begin{equation}\label{eq:normalisation}
u=\frac{D-B}{W-B},\qquad W>B.
\end{equation}
Si $D=6144$, $B=2048$ et $W=16383$, alors $u\approx0{,}28574$.
Changer arbitrairement $W$ modifie l'échelle de tous les pixels.
Changer $B$ affecte proportionnellement davantage les ombres.

Le blanc opérationnel peut être inférieur au maximum théorique
$2^b-1$. Il dépend parfois du canal, de l'ISO ou du mode. Les marges
optical-black sont des zones masquées destinées à mesurer le noir,
mais toutes les bordures du tableau ne sont pas nécessairement masquées.
Une zone active prise pour du noir produit une correction erronée.

Au-dessus de la capacité de saturation, plusieurs éclairements
distincts donnent le même code. Aucune transformation réversible
du fichier ne peut les distinguer. Une reconstruction de hautes
lumières peut exploiter des canaux non saturés ou une hypothèse de
couleur ; elle n'est pas une récupération certaine de la mesure perdue.

\section{Échantillonnage spatial et résolution}
\subsection{Du champ continu à la grille}
Une image monochrome discrète est une fonction
$I:\{0,\ldots,W-1\}\times\{0,\ldots,H-1\}\to\mathbb R$.
Le pixel représente une mesure spatialement intégrée, pas un point
mathématique infiniment petit. Une grille $6000\times4000$ contient
24 millions de positions. Cela ne garantit ni 24 millions de détails
indépendants ni une résolution optique correspondante.

Une fréquence spatiale trop élevée peut se replier en fréquence
plus basse : c'est l'aliasing. En une dimension, une sinusoïde à
$0{,}7$ cycle par échantillon peut être confondue, sur la grille,
avec une fréquence de $0{,}3$. La limite de Nyquist est $0{,}5$ cycle
par échantillon pour un signal correctement limité en fréquence.
En photographie, les structures fines, tissus et grilles peuvent
produire du moiré ; le CFA ajoute sa propre structure d'échantillonnage.

\subsection{Redimensionner n'est pas seulement sélectionner des pixels}
Réduire une image par deux doit généralement être précédé d'un
passe-bas. Conserver un pixel sur deux sans filtrer peut transformer
une texture fine en motif faux. Agrandir reconstruit une fonction
interpolée : plus proche voisin, bilinéaire, bicubique ou autre.
Le choix change netteté, halos et aliasing ; il n'ajoute pas de
mesures physiques nouvelles.

Un filtre bilinéaire en deux dimensions combine quatre voisins :
\begin{equation}
I(x,y)=(1-a)(1-b)I_{00}+a(1-b)I_{10}
+(1-a)bI_{01}+abI_{11},
\end{equation}
où $a,b\in[0,1]$ sont les fractions de position. Si les valeurs sont
des couleurs encodées non linéaires, cette moyenne n'est pas la
même que la moyenne des lumières ; nous y reviendrons.

\subsection{DPI, PPI, dimensions et taille physique}
Une largeur de 6000 pixels imprimée à 300 pixels par pouce donne
$6000/300=20$ pouces, soit $50{,}8$ cm. Changer seulement la métadonnée
de résolution à 150 PPI ne crée aucun pixel : la dimension physique
demandée devient 101,6 cm. Le DPI d'une imprimante décrit ses points
matériels et n'est pas nécessairement égal au nombre de pixels d'image.

\section{Mosaïque colorée, canaux et profondeur}
\subsection{Photosite, CFA et pixel RGB}
Dans un capteur Bayer courant, chaque photosite mesure une réponse
spectrale principalement rouge, verte ou bleue. Une cellule RGGB est
\[
\begin{matrix}R&G\\G&B\end{matrix}.
\]
Ce sont des filtres larges et chevauchants, pas des monochromateurs
parfaits. Les deux verts occupent des positions différentes ; on
peut devoir préserver quatre classes RGGB dans la calibration.
Le dématriçage estime les composantes absentes à chaque position.
Le RAW mosaïqué de 24 millions de photosites contient ainsi environ
24 millions de mesures, et non 72 millions de mesures RGB indépendantes.

Des capteurs monochromes, des CFA différents, des capteurs à couches
et des modes multivues existent. Il ne faut jamais déduire le motif
de la seule extension du fichier. Le décalage du recadrage peut aussi
changer la phase du motif : enlever une colonne fait commencer une
ligne RG par GR.

\subsection{Profondeur par échantillon et profondeur par pixel}
Un entier non signé sur $b$ bits possède $2^b$ codes, de 0 à $2^b-1$.
Une image RGB 8 bits par composante utilise 24 bits par pixel avant
compression ; RGBA utilise 32 bits. Un RAW 14 bits mosaïqué n'est
donc pas « moins riche » qu'un JPEG 24 bits au sens où ces nombres
seraient directement comparables : ils décrivent des objets différents.

Une quantification uniforme normalisée est
\begin{equation}
Q_b(x)=\frac{\operatorname{round}((2^b-1)x)}{2^b-1},\quad x\in[0,1].
\end{equation}
Le pas est $\Delta=1/(2^b-1)$ et l'erreur d'arrondi est bornée par
$\Delta/2$. Sous des hypothèses statistiques supplémentaires,
sa variance est approximativement $\Delta^2/12$ ; cette erreur n'est
pas systématiquement un bruit aléatoire indépendant. Sur un dégradé
très lisse, elle peut produire des bandes. Le dithering ajoute un
signal contrôlé avant quantification pour réduire ces structures,
au prix d'un bruit visible ou mesurable.

\subsection{Profondeur nominale, utile et stockage}
Un signal 12 bits peut être rangé dans un entier 16 bits. Il faut
savoir s'il occupe les bits faibles, les bits forts ou une
représentation compactée. La profondeur de stockage ne démontre
pas le nombre de bits utiles du capteur. De même, convertir un
JPEG 8 bits en TIFF 16 bits ne recrée pas les valeurs intermédiaires
perdues : cela donne surtout davantage de marge pour les calculs suivants.

\resultat{Résolution spatiale, nombre de composantes, précision numérique,
dynamique physique, gamut et compression sont des dimensions distinctes
de la qualité. Aucun nombre unique ne les résume.}

\part{Comment les nombres deviennent des octets}
\section{Binaire, hexadécimal et ordre des octets}
\subsection{Un octet n'a pas de type intrinsèque}
Un octet possède huit bits, donc 256 configurations. La suite
\texttt{01000001} vaut 65 en entier non signé et représente aussi
la lettre A en ASCII. Le type est une règle d'interprétation, pas
une propriété magique de ces huit bits. En hexadécimal, un chiffre
représente quatre bits : \texttt{41} désigne le même octet.

Pour des bits $b_i\in\{0,1\}$,
\begin{equation}
n=\sum_{i=0}^{k-1}b_i2^i.
\end{equation}
Le bit d'indice 0 est le bit de poids faible. Le nombre 291 vaut
\texttt{0x0123} et $1\times256+2\times16+3$.
Un affichage hexadécimal montre généralement une adresse à gauche,
les octets au centre et une interprétation texte à droite. Le texte
illisible ne signifie pas que le fichier est chiffré : des pixels
et des nombres binaires ne sont pas du texte.

\subsection{Little-endian et big-endian}
Pour l'entier 16 bits \texttt{0x1234}, big-endian écrit
\texttt{12 34}, little-endian écrit \texttt{34 12}.
L'ordre des octets n'est pas la même question que l'ordre des bits
dans un flux compacté. Un conteneur et son codec interne peuvent
avoir des conventions différentes.

Lire deux octets $a,b$ en big-endian donne $256a+b$ ; en
little-endian, $a+256b$. Interpréter \texttt{01 00} avec le mauvais
ordre transforme 1 en 256. Sur une image, cela produit souvent
une luminosité aberrante ou des motifs périodiques.

\subsection{Entiers signés et flottants}
En complément à deux, un entier signé de $b$ bits va de
$-2^{b-1}$ à $2^{b-1}-1$. Un résidu de compression peut être
négatif même si l'échantillon final est non signé.

Un flottant binaire normalisé peut être compris comme
\begin{equation}
(-1)^s\,2^{e-\mathrm{biais}}\left(1+\frac{m}{2^p}\right).
\end{equation}
Il distribue une précision approximativement relative sur une
large plage. Le type binary32 utilise 1 bit de signe, 8 d'exposant
et 23 de fraction ; binary16 utilise 1, 5 et 10. Il existe aussi
zéros, subnormaux, infinis et NaN. Les valeurs très proches de zéro
ne suivent pas toutes la formule normalisée ci-dessus.

Un flottant 16 bits n'est donc pas un entier 16 bits : il peut
représenter des valeurs négatives et supérieures à 1, mais ne
possède pas 65 536 niveaux uniformément répartis entre 0 et 1.
Dans un moteur, les NaN et infinis doivent être détectés : une
seule valeur non finie peut contaminer histogrammes et statistiques.

\section{Organisation d'une image en mémoire}
\subsection{Entrelacé et planaire}
Une image RGB peut être rangée
\texttt{R0 G0 B0 R1 G1 B1}, ou en trois plans
\texttt{R0 R1 ...}, \texttt{G0 G1 ...}, \texttt{B0 B1 ...}.
La première organisation facilite certaines opérations par pixel ;
la seconde facilite certaines opérations par canal et certains codecs.
RGB et BGR sont aussi deux ordres différents. Une permutation R/B
est un problème d'organisation, pas une balance des blancs.

Pour une image entrelacée de $C$ canaux, avec $s$ octets par composante :
\begin{equation}\label{eq:address}
\operatorname{adresse}(x,y,c)=A_0+y\,\mathrm{stride}+xCs+cs.
\end{equation}
Le \emph{stride}, ou pitch, est le nombre d'octets séparant les
débuts de deux lignes. Il peut dépasser $WCs$ pour respecter un
alignement ou parce que l'image est une vue dans un plus grand tampon.

\subsection{Exemple de lignes rembourrées}
Une ligne de 3 pixels RGB 8 bits contient 9 octets utiles.
Si chaque ligne est alignée sur 4 octets, le stride peut être 12.
Les 3 derniers octets sont du remplissage, pas un quatrième pixel.
Pour $H=2$, la mémoire allouée vaut alors 24 octets, pas 18.
Une boucle qui suppose des lignes contiguës de 9 octets lira mal
la deuxième ligne.

Une région d'intérêt peut conserver le stride de son parent.
Son pointeur commence plus loin, mais sa largeur n'explique pas
toute la distance entre ses lignes. En C++/OpenCV, ne présumer
la continuité d'un \texttt{cv::Mat} qu'après vérification ; dans Qt,
utiliser le pas réel de l'image. Une vue NumPy peut également être
non contiguë et même posséder un stride négatif.

\subsection{Origine, axes et orientation}
En traitement d'image, l'origine est souvent en haut à gauche,
$x$ croît vers la droite et $y$ vers le bas. En géométrie
mathématique, $y$ croît souvent vers le haut. Une rotation positive
peut donc sembler inversée si l'on confond transformation du repère
et transformation du contenu.

Il faut distinguer la grille stockée de l'orientation d'affichage.
Une photo portrait peut contenir une grille paysage avec une
instruction de rotation. Les huit orientations Exif couvrent
rotations et réflexions~\cite{exif}. Pour un exemple de rotation
horaire d'une grille $W\times H$ :
\[
x'=H-1-y,\qquad y'=x.
\]
Les dimensions deviennent $H\times W$. Si le logiciel tourne les
pixels tout en conservant une instruction de rotation non remise
à jour, le lecteur suivant peut effectuer une seconde rotation.

\subsection{Alpha droit et alpha prémultiplié}
L'alpha exprime une couverture ou opacité conventionnelle, pas une
quatrième couleur ni forcément une transparence physique complète.
Avec alpha droit, la couleur $C$ est stockée indépendamment de $a$.
Avec alpha prémultiplié, le tampon stocke $P=aC$.

Pour deux couches, avant-plan $f$ et arrière-plan $b$, en prémultiplié :
\begin{equation}
P_{\mathrm{out}}=P_f+(1-a_f)P_b,\qquad
a_{\mathrm{out}}=a_f+(1-a_f)a_b.
\end{equation}
Si l'on veut une couleur droite, diviser $P_{\mathrm{out}}$ par
$a_{\mathrm{out}}$ quand il est non nul. Les calculs de mélange
lumineux doivent se faire dans un espace linéaire approprié.
Dé-prémultiplier quand $a$ est presque nul amplifie les erreurs.
Des franges sombres autour des objets proviennent fréquemment
d'un mauvais contrat d'alpha ou d'une interpolation dans le
mauvais domaine.

\section{Coût mémoire et transfert GPU}
Pour $6000\times4000$ pixels :
\begin{center}
\begin{tabular}{lrr}
\toprule Représentation & Octets/pixel & Taille sans padding \\
\midrule Mosaïque stockée en uint16 & 2 & 48 MB \\
RGB uint16 & 6 & 144 MB \\
RGBA float32 & 16 & 384 MB \\
Deux tampons RGBA float32 & 32 & 768 MB \\
\bottomrule
\end{tabular}
\end{center}
Ici MB signifie $10^6$ octets ; 384 MB représentent environ
366,2 MiB, avec $1\,\mathrm{MiB}=2^{20}$ octets. Il faut ajouter
les intermédiaires, pyramides, masques, textures GPU et caches.
Le petit fichier JPEG de quelques mégaoctets peut ainsi produire
des centaines de mégaoctets en mémoire.

Un GPU ne lit pas nécessairement le TIFF ou le CR3 directement.
La chaîne usuelle comprend lecture, décompression CPU ou spécialisée,
allocation d'un tampon, transfert, traitement, puis éventuellement
retour CPU. Le coût total est
\[
T=T_{\text{lecture}}+T_{\text{décompression}}+
T_{\text{transferts}}+T_{\text{calcul}}+T_{\text{attentes}}.
\]
Un filtre très rapide ne compense pas toujours les transferts.
Mesurer chaque étage avant de conclure que « le GPU accélérera
l'ouverture ». Les formats de texture de rendu ne sont pas
nécessairement les formats photographiques de stockage.

\section{Un fichier est un conteneur structuré}
\subsection{Signature, en-tête, répertoires et charges utiles}
Une architecture générale peut être représentée ainsi :
\begin{center}
\begin{tikzpicture}[node distance=.45cm,
box/.style={draw=iprblue,rounded corners,text width=11cm,align=center,minimum height=.8cm}]
\node[box] (s) {Signature et version : quel contrat de lecture ?};
\node[box,below=of s] (h) {En-têtes et index : dimensions, types, offsets, compression};
\node[box,below=of h] (m) {Métadonnées et calibration : Exif, profils, orientation};
\node[box,below=of m] (d) {Charges utiles : pixels, coefficients, aperçus, masques};
\node[box,below=of d] (c) {Contrôles éventuels : longueurs, CRC, sommes de contrôle};
\draw[-{Latex}] (s)--(h);\draw[-{Latex}] (h)--(m);
\draw[-{Latex}] (m)--(d);\draw[-{Latex}] (d)--(c);
\end{tikzpicture}
\end{center}
Ce schéma est conceptuel : les sections peuvent être dispersées
dans le fichier, imbriquées ou répétées. Un TIFF utilise des
répertoires et offsets ; un PNG des chunks successifs ; un fichier
fondé sur ISO BMFF des boîtes imbriquées. Il n'existe pas un
en-tête universel de toutes les images.

\subsection{Codec, conteneur et extension}
Le conteneur organise les blocs ; le codec explique comment
reconstruire les échantillons de la charge utile. Une même
extension peut accepter plusieurs codecs. TIFF peut stocker des
valeurs non compressées ou différents flux compressés. CR3
désigne une famille de conteneurs Canon, non un simple tableau
RGB de convention fixe.

Renommer \texttt{photo.jpg} en \texttt{photo.dng} ne convertit rien.
Le type MIME, l'extension et la signature sont des indices à
croiser. Une signature valide prouve seulement que quelques
octets ressemblent à l'entrée d'un format, pas que tous ses blocs
sont corrects ou que son image se décodera.

\subsection{Offsets, longueurs et débordements}
Un offset est une distance en octets depuis une origine définie.
Cette origine peut être le début du fichier, celui d'un bloc ou
celui d'une structure interne. Une longueur peut compter les
données seules ou inclure l'en-tête. Ne pas deviner ces conventions.

Si un bloc commence à $o$ et contient $n$ octets dans un fichier
de taille $S$, vérifier $0\le o\le S$ puis $0\le n\le S-o$.
Cette seconde forme évite le débordement de $o+n$ dans un entier
de taille finie. De même, valider $W$, $H$, les canaux et le produit
avant d'allouer $WHCs$. Un fichier malveillant peut déclarer une
image gigantesque tout en ne contenant que quelques octets.

\subsection{CRC, hash et identité}
Un CRC détecte efficacement certaines corruptions accidentelles.
Un hash cryptographique comme SHA-256 sert à comparer des contenus
et à contrôler la provenance ; il ne prouve pas à lui seul
l'authenticité d'un auteur. Une signature cryptographique ajoute
une identité et une chaîne de confiance à vérifier.

Deux fichiers de hashes différents peuvent décoder les mêmes
pixels : ordre des métadonnées, compression et horodatages
diffèrent. À l'inverse, deux fichiers de même nom peuvent avoir
des contenus distincts. Il faut distinguer identité des octets,
égalité des pixels décodés et équivalence du rendu colorimétrique.

\resultat{Le fichier peut être compact, dispersé et compressé ; la
mémoire peut être large, alignée et flottante. Lire un fichier n'est
jamais simplement copier ses octets dans une texture.}

\part{Pourquoi et comment compresser une image}
\section{Réversibilité et théorie de l'information}
\subsection{Trois pertes à ne pas confondre}
La conversion analogique-numérique quantifie une mesure physique.
Un codec avec pertes peut ensuite supprimer de l'information
numérique. Enfin, un rendu peut écrêter le gamut ou les hautes
lumières. Une compression \emph{sans perte} garantit seulement
la reconstruction exacte des données auxquelles elle s'applique,
pas la préservation de toute la scène d'origine.

Si un encodeur $E$ et un décodeur $D$ sont sans perte sur un domaine
$\mathcal X$, on exige $D(E(x))=x$ pour tout $x\in\mathcal X$.
Une transformation inversible est une bijection sur ce domaine.
L'arrondi de 12,4 et de 12,49 vers 12 est au contraire plusieurs-à-un :
l'inverse exact est impossible sans information supplémentaire.

\subsection{Entropie : le coût moyen d'une surprise}
Si un symbole $i$ apparaît avec probabilité $p_i$, l'entropie est
\begin{equation}
H=-\sum_i p_i\log_2p_i.
\end{equation}
Un symbole certain apporte zéro bit de surprise ; deux symboles
équiprobables apportent un bit. Cette lecture probabiliste de
l'information vient de Shannon~\cite{shannon}. Elle ne mesure
ni la beauté de l'image ni sa valeur scientifique.

Exemple original : les probabilités $(1/2,1/4,1/8,1/8)$ donnent
$H=1{,}75$ bit par symbole. Un stockage fixe à deux bits consomme
deux bits par symbole, tandis qu'un code préfixe adapté peut
atteindre ici exactement 1,75 bit en moyenne. Le fichier réel
ajoute une description du code, des en-têtes et des limites de blocs.

Il est impossible de raccourcir sans perte \emph{tous} les fichiers :
il y a $2^n$ suites de $n$ bits et seulement $2^n-1$ suites de
longueur strictement inférieure à $n$. Une compression efficace
exploite donc des régularités de la famille d'images rencontrée.
Un bruit aléatoire indépendant se compresse souvent moins bien
qu'une plage uniforme, même si les deux ont les mêmes dimensions.

\section{Prédiction, différences et codes entropiques}
\subsection{Prédire pour rendre les données plus faciles à coder}
Si une ligne vaut $(100,102,101,103)$, conserver le premier nombre
et les différences donne $(100,2,-1,2)$. La transformation est
réversible : les sommes cumulées restituent la ligne.
Le nombre de valeurs n'a pas changé, mais les résidus se
concentrent autour de zéro et peuvent coûter moins de bits.

Avec un prédicteur $P_i$ calculé à partir de valeurs déjà connues :
\begin{equation}
r_i=x_i-P_i,\qquad x_i=P_i+r_i.
\end{equation}
Le décodeur doit employer exactement le même prédicteur et le
même ordre de parcours. Un prédicteur adaptatif doit préciser
son état initial et ses règles de mise à jour.

Une soustraction peut aussi être effectuée modulo 256 sur des
octets. Par exemple $3-250\equiv9\pmod{256}$, puis
$250+9\equiv3\pmod{256}$. Le résidu stocké est 9, pas un entier
signé de valeur $-247$. Cette convention permet des filtres
réversibles sans changer la largeur des octets.

\subsection{RLE : compter les répétitions}
La suite \texttt{AAAAABBBCC} peut être représentée par
$(5,A),(3,B),(2,C)$. Sur une photographie bruitée, les suites
constantes sont courtes ; ce codage peut même grossir les données.
Des marqueurs doivent distinguer une répétition d'une séquence
littérale. Le principe est simple, ses conventions binaires
exactes appartiennent au format qui l'utilise.

\subsection{Huffman : des mots courts pour les symboles fréquents}
Dans notre exemple de probabilités, on peut choisir
$A\mapsto0$, $B\mapsto10$, $C\mapsto110$, $D\mapsto111$.
Aucun mot n'est le début d'un autre : le code est préfixe.
La suite $ABCD$ devient \texttt{010110111} et se relit sans
séparateurs en parcourant un arbre. La longueur moyenne vaut
$0{,}5\times1+0{,}25\times2+0{,}125\times3+0{,}125\times3=1{,}75$.

Huffman code des symboles discrets avec des longueurs entières.
Un codage arithmétique ou par intervalles représente une suite
entière par une sous-plage, permettant une adaptation plus fine
aux probabilités. Les implémentations doivent définir les
arrondis et la renormalisation pour rester synchronisées.

\subsection{LZ77 et DEFLATE}
Un dictionnaire glissant remplace une séquence répétée par une
distance en arrière et une longueur. Il ne cherche pas des
objets visuels : il cherche des octets déjà rencontrés.
DEFLATE combine des références de ce type et un codage Huffman ;
son format est défini par la RFC~1951~\cite{deflate}.
Zlib ajoute un encadrement et un contrôle d'intégrité décrits
par la RFC~1950~\cite{zlib}. Zlib, DEFLATE et ZIP sont donc
trois notions apparentées, mais non interchangeables.

\subsection{Rice/Golomb : encoder des petits résidus}
Pour un entier non négatif $u$ et un paramètre $k$, un code de
Rice sépare quotient $q=\lfloor u/2^k\rfloor$ et reste
$r=u\bmod2^k$. Le quotient est écrit en unaire, le reste sur
$k$ bits. Avec la convention « $q$ uns suivis d'un zéro »,
$u=13$, $k=2$ donne $q=3$, $r=1$, donc \texttt{1110 01}.

Pour transformer un résidu signé en entier non négatif, on peut
utiliser $u=2r$ si $r\ge0$, sinon $u=-2r-1$. Les valeurs
$0,-1,1,-2,2$ deviennent $0,1,2,3,4$. Le choix de $k$ et les
conventions varient selon le codec ; cet exemple n'est pas un
décodeur Canon, Sony ou Nikon.

\section{Transformées : changer de coordonnées}
\subsection{Pourquoi passer aux fréquences ?}
Un bloc uniforme contient beaucoup de valeurs égales. Dans une
base adaptée, il peut être représenté par un coefficient moyen
et des détails nuls. Une transformée inversible ne compresse
pas nécessairement à elle seule ; elle concentre l'information
pour une quantification et un codage plus efficaces.

Pour une transformation orthonormale $A$, $c=Ax$ et
$x=A^{\mathsf T}c$. L'énergie quadratique se conserve :
$\|c\|_2^2=\|x\|_2^2$. Les erreurs dans les coefficients se
traduisent alors directement en erreurs quadratiques dans
l'image, sous réserve d'arrondis de calcul.

\subsection{DCT 2D et exemple JPEG}
Une DCT orthonormale d'un bloc $8\times8$ centré est
\begin{equation}
F(u,v)=\frac14\alpha(u)\alpha(v)
\sum_{x=0}^{7}\sum_{y=0}^{7}f(x,y)
\cos\!\frac{(2x+1)u\pi}{16}
\cos\!\frac{(2y+1)v\pi}{16},
\end{equation}
avec $\alpha(0)=1/\sqrt2$, $\alpha(k)=1$ sinon.
Le coefficient $(0,0)$ est le DC ; les autres décrivent des
variations spatiales AC. Dans un exemple 8 bits classique,
on soustrait 128 avant la transformée.

Si tous les pixels valent 140, le bloc centré vaut 12 ;
$F(0,0)=96$ et les autres coefficients sont nuls en calcul exact.
L'inverse retrouve 12 partout, puis 140 après réaddition de 128.
Cette construction mathématique explique la concentration de
l'énergie ; le cadre normatif JPEG est T.81~\cite{jpeg}.

\subsection{Quantification : l'étape irréversible}
On calcule pour chaque coefficient
\begin{equation}
\widehat F(u,v)=\operatorname{round}\!\frac{F(u,v)}{Q(u,v)},
\qquad F_{\mathrm{reconstruit}}=Q\widehat F.
\end{equation}
Avec $F=23$ et $Q=10$, on stocke 2 et on reconstruit 20.
Les valeurs de 15 à moins de 25, selon la convention d'arrondi,
peuvent partager ce code ; l'information est perdue.
Les hautes fréquences sont souvent quantifiées plus fortement.
Des détails faibles deviennent nuls, ce qui favorise ensuite
les suites de zéros et les codes entropiques.

La « qualité JPEG 90 » est un paramètre d'encodeur, pas une
unité normalisée garantissant la même table entre bibliothèques.
Réexporter plusieurs fois avec pertes peut accumuler des erreurs,
même si certaines configurations se stabilisent partiellement.

\subsection{Zigzag, DC différentiel et sous-échantillonnage}
Un parcours zigzag regroupe approximativement les coefficients
des basses vers les hautes fréquences. Les zéros consécutifs
peuvent être codés efficacement. Le DC est souvent prédit à
partir du bloc précédent ; les marqueurs de redémarrage
réinitialisent des états définis par le codec.

Avant le codage, une image peut être convertie en une composante
de luma et deux de chroma. Le format 4:2:0 usuel conserve une
grille de luma et deux grilles de chroma à demi-largeur et
demi-hauteur. Pour un rectangle de dimensions paires, cela
représente environ $1{,}5WH$ échantillons contre $3WH$ en 4:4:4.
L'emplacement exact des échantillons de chroma et les règles de
reconstruction comptent. Les textes colorés et contours
saturés peuvent souffrir davantage que les surfaces photographiques.

\subsection{Ondelettes et lifting}
Les ondelettes décomposent l'image par échelles plutôt qu'en
blocs DCT indépendants. Une étape de Haar produit, pour deux
valeurs, $(a,d)=((x_0+x_1)/\sqrt2,(x_0-x_1)/\sqrt2)$.
Appliquée sur lignes puis colonnes, elle donne un fond et
trois orientations de détails. Le cours 10 développe cette
construction et son lien avec l'estimation du bruit.

Le lifting construit des transformées par prédiction et mise
à jour. Avec des règles d'arrondi entières symétriques,
certaines variantes sont exactement réversibles. Quantifier
leurs sous-bandes produit en revanche un codec avec pertes.
Dire « ondelettes » ne suffit donc pas à conclure « sans perte ».

\section{Ce qu'un décodeur doit reconstruire}
Un décodeur JPEG à DCT ne lit pas directement les pixels : il
parse les marqueurs, reconstruit les tables, décode les
symboles, remet les coefficients dans les blocs, déquantifie,
effectue la DCT inverse, reconstruit les composantes et applique
les conversions demandées. Une décompression prédictive suit
une chaîne différente. Les détails normatifs fixent la
précision, les modes et les conventions.

\resultat{Décompresser exactement un fichier avec pertes ne restaure
pas les valeurs originales : cela restaure exactement, ou suivant
la précision autorisée du décodeur, la représentation déjà appauvrie
que l'encodeur a choisie.}

\part{Autopsie de formats représentatifs}
\section{Un format simple pour commencer : PPM}
Le format PPM binaire de Netpbm, variante P6, décrit une largeur,
une hauteur, une valeur maximale puis des composantes RGB~\cite{netpbm}.
Pour une image originale de deux pixels, rouge puis vert :
\begin{verbatim}
P6
2 1
255
[FF 00 00] [00 FF 00]
\end{verbatim}
Les crochets sont une annotation, pas des caractères à écrire.
L'en-tête est du texte, la charge utile est binaire. Une lecture
par lignes de texte appliquée aux pixels est donc incorrecte.
Avec ces seuls pixels, il n'existe ni miniature ni masque ni
profil ICC embarqué. Il faut connaître les conventions de
couleur et les limites de l'outil ; un tableau RGB brut est
beaucoup moins auto-descriptif qu'un format photographique riche.
La spécification PPM nominale décrit des composantes non linéaires
selon BT.709, pas des intensités linéaires. Des variantes rencontrées
en pratique utilisent pourtant sRGB ou du RGB linéaire : il ne faut
pas déduire une colorimétrie certaine de la seule extension
\texttt{.ppm}. Nos deux pixels extrêmes ne permettent pas de
distinguer ces fonctions de transfert.

Cet exemple permet d'isoler la différence entre matrice et
fichier : le tableau fait six octets, le fichier en fait
davantage à cause de son en-tête. Le script compagnon produit
ce fichier sans utiliser de bibliothèque d'image.

\section{TIFF : un graphe de répertoires et de données}
\subsection{Le contrat structurel}
TIFF classique commence par un ordre des octets, l'identifiant
42 et l'offset du premier IFD. Un IFD contient un nombre
d'entrées, des entrées de 12 octets et l'offset de l'IFD suivant.
Chaque entrée décrit un tag, un type, un nombre de valeurs et
une valeur courte ou un offset. Dimensions, profondeur,
compression, organisation et emplacement des bandes sont
des tags, pas un bloc universel de pixels placé à une adresse
constante~\cite{tiff}.

\subsection{En-tête original commenté octet par octet}
Un petit TIFF little-endian peut commencer ainsi :
\begin{verbatim}
Offset  Octets                  Interpretation
0       49 49                   "II", little-endian
2       2A 00                   entier 16 bits : 42
4       08 00 00 00             premier IFD a l'offset 8
8       09 00                   neuf entrees dans cet exemple
10      00 01 04 00 01 00 ...   debut du tag ImageWidth
\end{verbatim}
\texttt{00 01} vaut ici \texttt{0x0100}=256 : c'est
l'identifiant du tag ImageWidth, pas la largeur. Le type
\texttt{04 00} correspond à LONG ; le nombre de valeurs est
ensuite un entier sur quatre octets. Dans l'exemple du TP,
la valeur finale de cette entrée sera la largeur 2.

Le nombre 42 ne représente donc ni une taille de fichier ni
un nombre de pixels. Une entrée « largeur = 2 » doit être
distinguée du tag qui indique où chercher les pixels.

\subsection{Exemple de tags utiles}
\begin{longtable}{p{4.3cm}p{2cm}p{8.2cm}}
\toprule Nom & Identifiant & Rôle conceptuel \\ \midrule
ImageWidth / ImageLength & 256 / 257 & Dimensions de l'image concernée par cet IFD. \\
BitsPerSample & 258 & Profondeur par composante ; peut être une liste. \\
Compression & 259 & Méthode de reconstruction de la charge utile. \\
Photometric\-Interpretation & 262 & Sens des composantes : gris, RGB, palette, etc. \\
StripOffsets / StripByteCounts & 273 / 279 & Débuts et tailles des bandes. \\
Orientation & 274 & Transformation d'affichage. \\
SamplesPerPixel & 277 & Nombre de composantes à chaque position. \\
RowsPerStrip & 278 & Découpage vertical en bandes. \\
PlanarConfiguration & 284 & Composantes entrelacées ou séparées. \\
SampleFormat & 339 & Interprétation numérique : entier, flottant, etc. \\
\bottomrule
\end{longtable}
Le tableau sert de carte de lecture, pas de liste suffisante
pour rendre tout TIFF valide. Les valeurs par défaut, les
contraintes par mode et les extensions doivent être vérifiées
dans les documents normatifs.

\subsection{Bandes, tuiles, pages et sous-images}
Une image en bandes peut se lire par portions horizontales.
Une image en tuiles permet un accès spatial par blocs, utile
pour les très grandes images. Plusieurs IFD peuvent décrire
des pages ou des résolutions ; des SubIFD peuvent pointer vers
des images associées. L'ordre physique dans le fichier ne
définit pas forcément l'ordre logique des images.

Une compression sans perte conserve les échantillons stockés,
mais un TIFF peut aussi contenir une charge utile avec pertes.
Un TIFF flottant peut porter une image linéaire HDR ; un TIFF
entier 16 bits sRGB peut être une image SDR déjà rendue. Il
faut lire les tags et le profil au lieu de conclure à partir
de l'extension.

\subsection{BigTIFF et limite des offsets}
TIFF classique utilise de nombreux champs 32 bits, ce qui
limite l'adressage. BigTIFF introduit un en-tête et des
structures avec offsets plus larges~\cite{bigtiff}. Ce n'est
pas « un TIFF classique avec un plus gros fichier » : un
lecteur doit reconnaître la variante avant de lire ses IFD.
Un dépassement de capacité peut sinon produire des offsets
tronqués et un export apparemment réussi, mais inutilisable.

\section{PNG : chunks, prédiction et DEFLATE}
\subsection{Structure normative résumée}
PNG utilise une signature de huit octets puis des chunks.
Chaque chunk comporte longueur big-endian, type, données et
CRC sur type plus données. IHDR décrit notamment dimensions,
profondeur et type de couleur ; IDAT porte le flux zlib ; IEND
termine le fichier. Les IDAT concaténés forment un flux unique,
leurs frontières ne sont pas des frontières de lignes ou de
blocs DEFLATE. Les lignes décompressées portent un octet de
filtre à inverser. PNG prend en charge plusieurs types de
couleur et des profondeurs autorisées selon le type~\cite{png}.

\subsection{Un en-tête de deux pixels fabriqué pour le TP}
La signature en hexadécimal est :
\begin{verbatim}
89 50 4E 47 0D 0A 1A 0A
\end{verbatim}
Pour notre exemple RGB 8 bits de largeur 2 et hauteur 1,
les treize octets de données IHDR sont :
\begin{verbatim}
00 00 00 02   00 00 00 01   08 02 00 00 00
\end{verbatim}
La longueur du chunk vaut 13 ; son coût total est
$4+4+13+4=25$ octets. L'image de six octets utiles prend
donc nécessairement davantage de place que six octets dans
ce format. Le script calcule le CRC plutôt que d'en recopier
une constante opaque.

\subsection{Comprendre les filtres sans confondre octets et pixels}
Appelons $x$ l'octet courant, $a$ l'octet correspondant à gauche,
$b$ celui au-dessus et $c$ celui en haut à gauche. Un filtre
stocke $(x-P)\bmod256$ et le décodeur récupère
$(r+P)\bmod256$. Le voisin gauche est éloigné d'un pixel
complet en octets : pour RGB 8 bits, il est à trois octets,
pas forcément à l'octet précédent.

Exemple original : sur une ligne monochrome $(10,12,15)$,
la prédiction par le voisin gauche donne $(10,2,3)$.
Le premier voisin absent vaut zéro dans ce modèle. Une
prévision par la ligne précédente est efficace sur des
zones verticalement régulières. Une prédiction combinant
gauche, dessus et diagonale peut mieux suivre les gradients.
Le filtre ne supprime pas les différences : il les réexprime.

Les cinq types usuels sont None, Sub, Up, Average et Paeth.
Le TP implémente leur inversion sur des octets et vérifie
leur réversibilité sur des données synthétiques. Il ne
prend volontairement pas en charge toute la norme.

\subsection{Sans perte ne signifie pas identique en octets}
Deux encodeurs peuvent choisir des filtres, niveaux de
compression et découpages de chunks différents. Les fichiers
auront des SHA-256 distincts tout en redonnant le même RGB.
Une suppression du profil ou une conversion de couleur avant
l'encodage peut modifier le rendu, même si le codec PNG reste
sans perte sur les pixels qu'il a reçus.

\section{JPEG : marqueurs et flux de coefficients}
\subsection{Fichier et procédé de compression}
JPEG est une famille de procédés, pas exclusivement la variante
baseline avec pertes la plus courante. Les fichiers usuels
associent ces procédés à des conventions comme JFIF ou Exif.
On rencontre notamment SOI, DQT, SOF, DHT, SOS et EOI pour
l'entrée, les tables, les caractéristiques, les scans et la fin.
Les modes progressifs transmettent les coefficients en
plusieurs scans ; des procédés prédictifs sans perte existent
aussi dans T.81~\cite{jpeg}.

\subsection{Schéma de lecture d'un JPEG usuel}
\begin{verbatim}
SOI
  APP... : description JFIF, Exif, ICC, XMP selon conventions
  DQT    : tables de quantification
  SOF... : dimensions, composantes, facteurs d'echantillonnage
  DHT    : tables Huffman si ce mode les utilise
  SOS    : description du scan
  ...    : donnees entropiques du scan
  ...    : scans ou segments supplementaires selon le mode
EOI
\end{verbatim}
Ce schéma ne fixe pas un ordre unique pour tous les modes.
Dans les données entropiques, des conventions d'échappement
et de redémarrage évitent de prendre des octets de données
pour des marqueurs. Chercher simplement \texttt{FF D9}
dans n'importe quelle sous-structure n'est pas un parseur.

Une miniature JPEG peut elle-même être contenue dans un
JPEG ou un RAW plus grand. Il faut distinguer le flux principal
de l'aperçu. Un logiciel peut réussir à afficher l'aperçu
tout en échouant à décoder la charge utile principale.

\section{DNG : TIFF plus une description du négatif numérique}
DNG étend TIFF avec des informations de calibration et de
développement : CFA, zones actives, niveaux, matrices et
autres tags de pipeline. La spécification distingue notamment
des données mosaïquées et des données linéaires déjà
dématricées ; plusieurs compressions sont possibles selon
les versions et modes. Les tags d'un profil DCP décrivent
une transformation ou un rendu, pas les pixels de la
photographie~\cite{dng}.

Un DNG valide doit être cohérent comme conteneur \emph{et}
comme négatif numérique. Écrire trois canaux RGB dans un TIFF
puis ajouter une extension DNG ne recrée pas la mosaïque
originale ni toutes ses possibilités de développement.
Un DNG linéaire peut être légitime, mais il doit déclarer
correctement sa nature et ses calibrations. Si des effets
sont déjà incorporés aux pixels, il ne faut pas décrire
ces pixels comme s'ils étaient encore les mesures brutes.

\section{HEIF, HEIC, AVIF et conteneurs à boîtes}
HEIF organise des éléments d'image et leurs relations dans
une structure fondée sur ISO BMFF. HEIC est couramment
associé à des images codées en HEVC ; AVIF associe une image
AV1 à une organisation HEIF. La spécification AVIF décrit
notamment propriétés de couleur et relations d'éléments~\cite{avif}.

Une boîte peut commencer par une taille et un type de quatre
caractères, puis contenir des sous-boîtes ou une charge utile.
Des associations peuvent relier une image à une miniature,
un alpha ou d'autres données. Un parser doit savoir quelle
structure est un conteneur avant de descendre récursivement :
tous les octets d'une boîte ne sont pas d'autres boîtes.

L'existence de modes HDR, alpha ou haute précision ne signifie
pas que tout lecteur les gère. Vérifier les propriétés exactes
du fichier et la compatibilité de la chaîne de décodage.
Un codec moderne n'est pas automatiquement un espace couleur.

\section{OpenEXR et images de calcul scientifique ou de rendu}
OpenEXR permet des canaux nommés, notamment HALF, FLOAT et
UINT, des fenêtres de données, tuiles et parties multiples.
Certaines compressions sont sans perte, d'autres avec
pertes. Des données profondes peuvent contenir plusieurs
échantillons par position ; les composantes couleur sont
habituellement linéaires et prémultipliées dans les usages
de compositing~\cite{openexr}.

Un canal Z peut représenter la profondeur et non une couleur.
Une composante supérieure à 1 peut porter une énergie
relative élevée sans être « invalide ». Elle ne devient
blanche sur un écran qu'après choix d'un rendu. L'architecture
d'un fichier image dépasse ainsi largement le triplet RGB.

\section{Panorama et choix selon l'usage}
\begin{longtable}{p{2.1cm}p{4.1cm}p{4.4cm}p{3.9cm}}
\toprule Famille & Atout fréquent & Attention & Usage typique \\ \midrule
JPEG usuel & Taille modeste, diffusion & Pertes et chroma possibles & Livraison de photos SDR. \\
PNG & Réversibilité des échantillons & Coût sur photographies bruitées & Graphiques, captures, alpha. \\
TIFF & Souplesse et haute précision & Variantes nombreuses & Échange et comparaison contrôlée. \\
RAW constructeur & Mesures et calibration & Support dépendant du modèle/mode & Développement photographique. \\
DNG & Description documentée du négatif & Cohérence et version à respecter & Échange/archivage RAW ou linéaire. \\
AVIF/\allowbreak HEIF & Codage moderne, éléments associés & Compatibilité variable & Diffusion et usages HDR compatibles. \\
OpenEXR & Flottants et canaux génériques & Pas un rendu photo universel & VFX, intermédiaires, calcul HDR. \\
\bottomrule
\end{longtable}
Il existe aussi BMP, GIF, WebP, JPEG~2000, JPEG~XL, PSD et
des formats médicaux ou géographiques. Leurs objectifs
diffèrent : animation, couches, séries, coordonnées,
profondeur ou accès aléatoire. Les principes étudiés ici
restent utiles pour les aborder sans supposer qu'ils suivent
tous le modèle JPEG ou TIFF.

\part{Des échantillons décodés à une photographie interprétée}
\section{Décodage, développement et rendu : trois contrats}
\subsection{Décodage du fichier}
Le décodage identifie le format, lit les structures, localise
les charges utiles et reconstruit les échantillons déclarés.
Pour un RAW mosaïqué, la sortie peut rester une seule grille
de codes par photosite. Pour un TIFF RGB, elle peut être un
tableau de composantes déjà rendues. Les deux sorties n'ont
pas la même signification, même si elles utilisent uint16.

\subsection{Développement du RAW}
Le développement utilise les mesures et les calibrations pour
produire une image interprétable : niveaux, corrections,
balance des blancs, reconstruction colorée, transformation
vers un espace de travail et éventuellement rendu de base.
LibRaw sépare notamment ouverture, dépaquetage et étapes de
traitement dans son API ; lire \texttt{open\_file} ou réussir
\texttt{unpack} n'est pas obtenir automatiquement un rendu
final~\cite{libraw}.

\subsection{Rendu et affichage}
Le rendu décide comment représenter une scène dans les
contraintes du support. Il comprend éventuellement courbes,
compression des hautes lumières, choix esthétiques, réduction
du gamut et conversion vers l'écran. Il n'existe pas un
seul rendu mathématiquement obligatoire pour toute scène.
En revanche, une interprétation incorrecte du noir, de la
matrice ou de la fonction de transfert est un défaut
technique, pas simplement un style artistique.

\resultat{Deux logiciels peuvent décoder les mêmes mesures puis
produire des rendus différents. Pour identifier un bug, il faut
comparer les étapes, pas seulement les images finales.}

\section{Un pipeline RAW conceptuel détaillé}
\begin{center}
\begin{tikzpicture}[node distance=.3cm,
box/.style={draw=iprblue,rounded corners,text width=12.5cm,align=center,minimum height=.66cm}]
\node[box] (f) {Fichier : structures, métadonnées, aperçus et données RAW};
\node[box,below=of f] (r) {Décompression et dépaquetage : grille native};
\node[box,below=of r] (n) {Niveaux, linéarisation éventuelle, pixels défectueux et corrections capteur};
\node[box,below=of n] (w) {Gains de balance des blancs et reconstruction des canaux};
\node[box,below=of w] (p) {Transformation caméra, profil, tables et rendu tonal de base};
\node[box,below=of p] (e) {Réglages, retouches, géométrie et traitement du détail};
\node[box,below=of e] (o) {Rendu de sortie, gamut, fonction de transfert, quantification et fichier};
\draw[-{Latex}] (f)--(r);\draw[-{Latex}] (r)--(n);\draw[-{Latex}] (n)--(w);
\draw[-{Latex}] (w)--(p);\draw[-{Latex}] (p)--(e);\draw[-{Latex}] (e)--(o);
\end{tikzpicture}
\end{center}
Cet ordre est pédagogique, pas universel. Certains débruitages
opèrent avant dématriçage ; certaines corrections optiques
après reconstruction ; des profils imposent un domaine
et un ordre précis. Toute permutation doit être justifiée.

\subsection{Dépaquetage et précision native}
Un codec peut reconstruire un flux de codes 12 ou 14 bits,
qu'il faut ensuite ranger dans un tampon. Le dépaquetage
traduit leur disposition binaire, sans décider de leur couleur.
Un ordre des bits incorrect peut être masqué sur des faibles
valeurs et devenir visible sur les plus fortes. Les tests
doivent inclure zéro, maximum et valeurs traversant les
frontières d'octets.

\subsection{Linéarisation, noir et blanc}
Certains fichiers décrivent une table de linéarisation ou
des niveaux qui ne se résument pas à deux constantes. Une
normalisation doit respecter les données déclarées. Dans
le cas simple, on utilise l'équation~\eqref{eq:normalisation}.
Si les noirs varient entre canaux, lignes ou colonnes,
écrire $B_c(x,y)$ plutôt qu'une constante universelle.

Les valeurs négatives après soustraction du noir peuvent
représenter des fluctuations normales autour de zéro.
Les couper immédiatement introduit un biais positif dans
les statistiques des ombres. À l'inverse, certaines étapes
comme des tables définies seulement sur $[0,1]$ exigent une
politique explicite de bornage. Les deux impératifs doivent
être conciliés, pas ignorés.

\subsection{Correction des défauts et optique}
Les photosites chauds ou morts peuvent être remplacés par
une estimation locale, éventuellement adaptée au CFA.
Une correction de vignettage multiplie spatialement les
valeurs ; elle amplifie aussi le bruit des zones sombres.
Une correction géométrique rééchantillonne la grille ; elle
change donc les valeurs et les correspondances pixel à pixel.
Une correction d'aberration chromatique peut déformer
différemment les canaux. Cela rend une MAE sans recalage
géométrique difficile à interpréter.

\subsection{Balance des blancs}
Un modèle diagonal courant est
\begin{equation}
\begin{pmatrix}R'\\G'\\B'\end{pmatrix}=
\begin{pmatrix}g_R&0&0\\0&g_G&0\\0&0&g_B\end{pmatrix}
\begin{pmatrix}R\\G\\B\end{pmatrix}.
\end{equation}
Les gains tentent de rendre neutre une référence éclairée
par une lumière donnée dans un contexte de calibration.
Multiplier les gains par une constante commune change
aussi la luminosité ; leur normalisation importe.
Un Kelvin affiché est une description de chromaticité ou
un paramètre d'interface, pas les trois gains eux-mêmes.
Deux moteurs peuvent afficher des Kelvin différents à
partir de gains proches, ou l'inverse, suivant leurs modèles.

Sous lumière mixte, une seule balance globale ne peut
neutraliser toutes les régions. Sous un projecteur à spectre
discontinu, « corriger la température » ne recrée pas les
longueurs d'onde absentes. La nuance fournit un degré
supplémentaire, mais pas une reconstruction spectrale complète.

\subsection{Dématriçage : une estimation, pas une lecture directe}
Sur une mosaïque RGGB, on peut estimer un vert à la position
d'un rouge par la moyenne des quatre voisins verts :
\begin{equation}
\widehat G(x,y)=\tfrac14[G(x-1,y)+G(x+1,y)+G(x,y-1)+G(x,y+1)].
\end{equation}
C'est une approximation bilinéaire simple. Elle lisse des
contours et peut produire des fausses couleurs. Des méthodes
adaptatives choisissent une direction ou exploitent la
corrélation des différences de couleur. Les méthodes
apprises proposent d'autres compromis, mais exigent des
données d'entraînement et une validation de généralisation.

La comparaison de deux dématriçages doit séparer détail,
bruit, moiré et rendu couleur. Plus net ne signifie pas
nécessairement plus fidèle : une arête peut être un artefact.

\section{La couleur n'est pas contenue dans les nombres seuls}
\subsection{Du spectre aux composantes}
Une réponse de caméra peut s'écrire
\begin{equation}
c_k=\int S(\lambda)\rho(\lambda)s_k(\lambda)\,d\lambda.
\end{equation}
Les fonctions $s_k$ sont les sensibilités spectrales effectives
des canaux. Les tristimulus d'un observateur standard
s'obtiennent avec d'autres fonctions de pondération, par
exemple $\overline x,\overline y,\overline z$ pour XYZ.
Trois nombres de caméra ne sont donc pas directement les
trois nombres d'un espace RGB d'affichage.

Deux spectres distincts peuvent donner les mêmes trois
réponses : c'est une forme de métamérisme. Une matrice
3$\times$3 peut approximer une transformation sur un domaine
donné, mais ne reconstitue pas tout le spectre. Une charte
photographiée sous un illuminant fournit une calibration
utile ; elle ne garantit pas la perfection sous tous les
éclairages et pour toutes les matières.

\subsection{Matrices et changement de base}
Dans un modèle linéaire :
\begin{equation}
\begin{pmatrix}X\\Y\\Z\end{pmatrix}=M
\begin{pmatrix}R_c\\G_c\\B_c\end{pmatrix}.
\end{equation}
Chaque colonne indique la contribution d'un canal caméra.
La direction de la matrice est essentielle : caméra vers
XYZ et XYZ vers caméra sont deux transformations inverses,
pas des noms interchangeables. Certaines calibrations
incluent des conventions de balance ou d'adaptation.
Appliquer deux fois une correction qui était déjà incorporée
produit une dérive.

Un espace RGB est défini par des primaires, un point blanc
et une convention de codage. Le passage entre primaires
linéaires peut être une matrice ; le passage entre codes
non linéaires exige d'abord l'inversion de leurs fonctions
de transfert. Une matrice appliquée directement à des
codes sRGB ne réalise généralement pas le bon changement
de base physique.

\subsection{Point blanc et adaptation chromatique}
D50 et D65 sont des blancs de référence, pas les températures
mesurées d'une photo particulière. L'adaptation chromatique
modélise une partie de l'adaptation visuelle lorsque l'on
change de blanc. Un modèle de type von Kries transforme
vers un espace de réponses, multiplie par des gains, puis
revient : $M^{-1}\operatorname{diag}(d)M$.
Bradford est un choix de matrice de ce type, développé dans
le cours 03. Confondre D50 et D65 peut déplacer les couleurs
même lorsque toutes les multiplications sont numériquement
correctes.

\section{Fonctions de transfert, gamma et HDR}
\subsection{Linéaire veut dire linéaire par rapport à quoi ?}
Dans un domaine linéaire par rapport à la lumière, doubler
un signal non saturé double la grandeur représentée. Une
valeur sRGB de 0,5 n'est pas la moitié de la lumière de 1.
La fonction d'encodage sRGB normalisée est~\cite{srgbimages}
\begin{equation}
E(L)=\begin{cases}
12{,}92L,&L\le0{,}0031308,\\
1{,}055L^{1/2{,}4}-0{,}055,&L>0{,}0031308.
\end{cases}
\end{equation}
Son inverse usuel est
\begin{equation}
L(E)=\begin{cases}
E/12{,}92,&E\le0{,}04045,\\
((E+0{,}055)/1{,}055)^{2{,}4},&E>0{,}04045.
\end{cases}
\end{equation}
Ainsi $E=0{,}5$ donne $L\approx0{,}21404$ et
$L=0{,}18$ donne $E\approx0{,}46136$. « Gamma 2,2 » est une
approximation fréquente, pas la fonction sRGB exacte.

\subsection{Pourquoi l'ordre des opérations compte}
La moyenne linéaire d'un noir et d'un blanc vaut $L=0{,}5$,
soit $E\approx0{,}73536$ après encodage. Faire la moyenne
des codes donne $E=0{,}5$, soit seulement $L\approx0{,}21404$.
On comprend ainsi pourquoi une interpolation dans le mauvais
domaine peut produire des transitions trop sombres.

Une correction d'exposition physique utilise
$L'=2^eL$, et non généralement $E'=2^eE$.
À $+1$ EV, un code sRGB 0,5 devient environ 0,68584
tant que le signal reste dans le domaine autorisé ; il ne
doit pas automatiquement devenir 1.

Quantifier avant l'encodage puis après lui n'est pas non
plus équivalent :
\[
Q(E(L))\ne Q(E(Q(L))).
\]
Une quantification linéaire prématurée peut supprimer des
faibles valeurs que l'encodage aurait réparties sur davantage
de codes. Ce problème a motivé des contrôles spécifiques
de sortie sRGB dans IPR ; ce n'est pas une propriété
inévitable de tout TIFF.

\subsection{OETF, EOTF et courbe artistique}
Une OETF traduit un signal optique vers un signal électrique
dans une chaîne donnée ; une EOTF décrit le comportement
d'un affichage à partir de ses codes. Une transformation
globale de la scène à l'écran peut aussi inclure une OOTF.
La courbe artistique ajuste le rendu tonal. Ce sont des
objets différents, même si plusieurs sont non linéaires.
Appliquer une courbe de profil ne dispense pas de décrire
l'encodage final ; appliquer deux encodages sRGB assombrit
ou éclaircit fortement le résultat selon la direction.

\subsection{HDR n'est pas synonyme de 16 bits}
HDR peut décrire une scène, un espace de calcul étendu ou
une diffusion pour un écran à grande luminance. Un TIFF
16 bits peut rester SDR ; un fichier à 10 bits correctement
décrit peut porter un rendu HDR. La recommandation
BT.2100 encadre notamment PQ et HLG~\cite{hdr}.

PQ relie les codes à une luminance d'affichage absolue selon
une fonction perceptuelle normalisée. HLG emploie une
construction différente, avec un comportement système lié
aux conditions d'affichage. Il faut conserver la fonction
de transfert, les primaires, la plage de codes et les
métadonnées pertinentes. Ouvrir les codes HDR comme sRGB
donne une image incohérente, même sans erreur de décompression.

Une \emph{gain map} peut décrire comment transformer une
base SDR vers un rendu à plus grande dynamique. Elle n'est
ni nécessairement une deuxième photographie indépendante,
ni une preuve que le capteur avait conservé toutes les
hautes lumières. Un moteur doit savoir s'il la prend en
charge ; l'ignorer peut laisser un rendu SDR utilisable,
mais différent du rendu voulu sur écran HDR.

\section{Profils ICC, DCP et LCP : des rôles distincts}
\subsection{ICC : décrire une relation colorimétrique}
Un profil ICC décrit une transformation entre un espace
de périphérique ou de travail et un espace de connexion
normalisé, suivant sa classe et sa structure. Des profils
peuvent utiliser matrices/courbes ou tables. Le PCS des
profils ICC classiques utilise notamment une référence D50 ;
les intentions de rendu guident le traitement des couleurs
non reproductibles~\cite{iccimages}.

\emph{Assigner} un profil change l'interprétation des nombres
sans forcément les changer. \emph{Convertir} cherche à
préserver l'apparence ou une relation colorimétrique en
calculant de nouveaux nombres. Assigner sRGB à des valeurs
Adobe RGB ne remplace donc pas une conversion Adobe RGB
vers sRGB. Un profil écran est destiné à la sortie du
système d'affichage, pas à devenir automatiquement le
profil incorporé à chaque photo exportée.

\subsection{DCP : caractériser et rendre des valeurs caméra}
Un DCP peut comporter matrices, illuminants, tables de
teinte/saturation, look, courbe et offset d'exposition.
Ces éléments n'ont pas tous un rôle purement colorimétrique :
certains participent au rendu de base. Leur domaine et leur
ordre sont essentiels~\cite{dng}. Une table décrite dans
des primaires de travail précises ne doit pas être évaluée
comme si elle était une table sRGB quelconque.

Une interpolation entre deux illuminants peut utiliser
des températures réciproques. Par exemple,
\begin{equation}
w=\operatorname{clip}\left(
\frac{1/T-1/T_1}{1/T_2-1/T_1},0,1\right),\qquad
M=(1-w)M_1+wM_2.
\end{equation}
C'est une règle de mélange d'un modèle, pas une interpolation
exacte de tous les spectres lumineux possibles.

Pour une courbe RGB, appliquer $T$ à chaque canal peut
changer la position relative du canal médian et donc la
teinte. Une alternative applique la courbe aux minimum et
maximum puis interpole le médian. DCamProf explique
l'intérêt de cette stabilisation et les interactions
avec les tables de rendu~\cite{dcamprofimages}. Elle ne
garantit ni saturation constante ni fidélité parfaite
après réduction du gamut. Le cours ne recommande pas
une substitution aveugle au moteur de production.

\subsection{LCP et corrections d'objectif}
Un profil d'objectif peut décrire distorsion, vignettage
et aberrations suivant focale, ouverture ou distance.
Son action peut être géométrique et photométrique. Il ne
remplace pas le profil couleur du capteur. Une correction
de distorsion déplace les coordonnées : pour des tests
pixel à pixel, son activation doit être harmonisée.

\subsection{Légalité et portabilité}
Une spécification de format ouverte et les données de
profils distribuées par un éditeur sont deux objets
différents. Le fait de pouvoir lire un profil ne donne
pas automatiquement le droit de le redistribuer avec
un logiciel. Les licences des bibliothèques, des profils,
des chartes et des jeux d'images doivent être examinées
séparément. Cette distinction n'est pas un avis juridique.

\section{Les métadonnées : plusieurs familles et plusieurs vérités}
\subsection{Inventaire utile}
\begin{longtable}{p{3.1cm}p{6.5cm}p{5cm}}
\toprule Famille & Exemples & Conséquence \\ \midrule
Structure & Dimensions, profondeur, compression, CFA & Nécessaire au décodage. \\
Prise de vue & Temps, ouverture, ISO, focale, date & Contexte, pas nécessairement une transformation. \\
Calibration & Noir, blanc, gains, matrices & Change l'interprétation du signal. \\
Présentation & Orientation, profil, miniature & Change le rendu ou sa géométrie. \\
Éditorial & Auteur, légende, mots-clés, droits & Organisation et diffusion. \\
Géographie & Latitude, longitude, altitude & Information sensible à protéger. \\
Développement & Exposition, courbe, paramètres de profil & Recette, parfois déjà incorporée aux pixels. \\
Gestion & Note, drapeau, libellé, identifiant & Peut vivre dans un catalogue ou sidecar. \\
\bottomrule
\end{longtable}

\subsection{Exif, MakerNotes, IPTC et XMP}
Exif fournit des conventions de métadonnées photographiques
et des structures associées aux images ; la page CIPA
recense les versions de la norme, dont Exif~3.1 publiée
en 2026~\cite{exif}. Un champ ISO peut figurer en plusieurs
endroits ou avec plusieurs interprétations historiques.
Il faut afficher les groupes de tags, pas seulement leurs
noms courts.

Les MakerNotes contiennent des données propres au
constructeur. Elles peuvent avoir leurs propres offsets,
versions et conventions. ExifTool documente les champs
qu'il sait interpréter ; cette documentation est une
source d'implémentation précieuse, pas une norme constructeur
universelle~\cite{exiftoolimages}. Déplacer un bloc sans
ajuster ses offsets peut casser sa lecture.

IPTC couvre notamment les usages éditoriaux ; XMP propose
un modèle extensible sérialisé couramment en RDF/XML,
embarqué ou externe~\cite{xmpimages,iptcimages}. Un catalogue
peut aussi garder des réglages sans les écrire dans le RAW.
« Je vois une étoile dans Lightroom » ne prouve donc pas
que cette étoile est présente dans le fichier transmis.

\subsection{Recette et pixels déjà développés}
Un TIFF exporté depuis Lightroom peut contenir des
métadonnées de développement ainsi que des pixels où le
développement est déjà effectué. Réappliquer les réglages
XMP sur ces pixels applique potentiellement l'effet deux
fois. Un moteur doit identifier la nature de l'entrée,
pas interpréter tout paramètre présent comme une instruction
de rendu à exécuter.

Un sidecar \texttt{.ipr} ou \texttt{.xmp} est une recette
non destructive seulement si l'original est préservé et
si les paramètres, versions et ressources nécessaires
restent disponibles. Pour reproduire un rendu, conserver
aussi le profil exact, la version moteur et les conventions.

\subsection{Données absentes, inconnues ou contradictoires}
Une valeur absente n'est pas automatiquement zéro.
« GPS absent » ne signifie pas latitude 0, longitude 0.
« Profil absent » ne prouve pas sRGB. Une température
calculée par un logiciel n'est pas forcément la valeur
inscrite par le boîtier. La stratégie correcte est de
conserver provenance, statut connu/inconnu et hypothèses
de repli, avec un avertissement si la comparaison dépend
de cette hypothèse.

\section{Pourquoi le même RAW change d'apparence selon le logiciel}
Les différences peuvent provenir du dépaquetage, du noir,
de la normalisation, de la WB, du dématriçage, des matrices,
des tables, de la courbe de base, de la récupération des
hautes lumières, des corrections optiques, du bruit,
du gamut mapping ou de l'affichage. Une miniature Canon
est un rendu déjà produit par Canon, pas une vue directe
des valeurs CFA. La comparaison à l'œil de cette miniature
et d'un export ne permet pas de localiser l'étape différente.

Le bon protocole compare d'abord les mêmes fichiers,
mêmes profils, mêmes réglages, mêmes fonctions de transfert
et même géométrie, puis isole les étages. Un profil nommé
« Neutral » n'est pas forcément une matrice sans courbe.
« Tous les curseurs à zéro » ne signifie pas « aucun
traitement ». Cette distinction est centrale pour la
validation de la base colorimétrique d'IPR.

\part{Canon, Sony et Nikon : principes communs, variantes réelles}
\section{Comment lire les informations propres à une marque}
La marque n'est pas une spécification de codec. Un modèle,
un firmware, un mode de rafale, un type d'obturateur, une
taille RAW et un réglage de compression peuvent changer
les données. Il faut constituer une matrice de support
\emph{modèle $\times$ mode $\times$ firmware}, pas une
simple liste d'extensions.

Une documentation de rétro-ingénierie est souvent la
meilleure explication publique de détails binaires,
mais elle doit être corroborée par des fichiers et des
implémentations. Une ancienne constante noire ou un
offset de MakerNote peut être valide sur un modèle et
faux sur son successeur. Un nouveau boîtier peut réutiliser
le conteneur tout en introduisant une nouvelle calibration.

\section{Canon : CR2, CR3, C-RAW et styles de rendu}
\subsection{CR2 : une organisation héritée de TIFF}
Les travaux de Laurent Clévy décrivent le CR2 comme une
structure TIFF enrichie, avec répertoires, métadonnées,
aperçus et données RAW. Ils expliquent aussi le décodage
JPEG prédictif sans perte et le réassemblage de tranches
rencontrés dans les fichiers étudiés~\cite{clevycr2}.
Ce JPEG sans perte n'est pas le JPEG à DCT avec pertes
d'une photographie livrée au client.

Il faut donc séparer trois opérations : trouver la charge
utile, décompresser ses échantillons, puis remettre les
échantillons dans leur géométrie capteur. Un tableau de
la bonne taille peut malgré tout avoir ses segments
permutés. Les modes réduits ou particuliers ne doivent
pas être traités comme une mosaïque pleine taille sans
vérifier leur organisation.

\subsection{CR3 : conteneur à boîtes et codec CRX}
La description publique CR3 de Clévy identifie notamment
\texttt{ftyp}, \texttt{moov}, des pistes, des données
\texttt{mdat}, des aperçus et des blocs de métadonnées
Canon. Les blocs CMT peuvent embarquer des structures
de type TIFF. La famille CRX comprend des variantes
prédictives et des variantes à ondelettes et quantification
pour C-RAW ; les détails observés sont documentés comme
rétro-ingénierie, non comme une norme officielle Canon~\cite{clevycr3}.

Le décodeur CRX de LibRaw contient lui-même des structures
de tuiles, plans, sous-bandes et paramètres de quantification~\cite{librawcrx}.
Le lecteur doit donc interpréter les descripteurs pour
savoir quel procédé appliquer. Déduire le mode du seul
suffixe CR3 est insuffisant. Le CR3 n'est pas non plus un
JPEG renommé : il peut seulement en contenir des aperçus.

\subsection{RAW contre C-RAW : ce que dit Canon et ce qu'on infère}
Canon présente C-RAW comme une variante de taille réduite
conservant la souplesse de développement RAW, et signale
que certaines remontées importantes de sous-exposition
peuvent révéler davantage de bruit. L'article constructeur
ne publie pas une description exhaustive de l'algorithme
binaire~\cite{canonraw}.

Les analyses et le code CRX apportent une description
plus technique de certains modes. Il ne faut pas convertir
un comportement observé sur une génération en garantie
pour toutes les futures versions. Pour tester un débruiteur
ou des ombres extrêmes, consigner explicitement RAW ou
C-RAW et éviter de mélanger les deux sans contrôle.

\subsection{Picture Style, Camera Neutral et Adobe Standard}
Canon distingue des styles de rendu tels que Standard,
Neutral et Faithful et décrit leurs objectifs visuels.
Le manuel du R6 Mark III distingue aussi force, finesse
et seuil de netteté~\cite{canonstyle}. Ces paramètres
participent à un rendu ; ils ne sont pas les coefficients
d'une simple matrice RGB à copier dans un décodeur.

Le style Neutre du boîtier, un DCP Adobe Camera Neutral
et un profil Adobe Standard sont trois objets différents.
Un profil Camera Matching cherche un rapprochement de
rendu, pas nécessairement une reproduction bit à bit du
moteur constructeur. Les noms n'impliquent pas non plus
une neutralité physique parfaite sous tous les éclairages.

Le JPEG intégré au RAW peut refléter le style choisi sur
le boîtier alors que le logiciel tiers développe la
mosaïque avec son propre profil par défaut. L'utilisateur
voit alors une transition entre l'aperçu Canon et le
rendu tiers, sans que les mesures RAW aient été changées.

\subsection{Exemple local IPR : calibrations à ne pas universaliser}
L'audit des six scènes Canon du projet a notamment
examiné une nouvelle version de ColorData, des zones
optical-black et la précision des faibles valeurs en
sortie. Une erreur de sélection de bordure peut biaiser
le noir ; une mauvaise interprétation du blanc change
la normalisation. La comparaison à Lightroom a aussi
isolé des différences de tables et de courbe.

Ces résultats datés sont dans
\href{../validation/analyses/analyse-base-colorimetrique.md}{le rapport évolutif de base colorimétrique}.
Les liens PDF relatifs dépendent du lecteur ; le chemin
depuis la racine est
\path{docs/validation/analyses/analyse-base-colorimetrique.md}.
Le cours retient la méthode, pas des valeurs à appliquer
à tous les RAW Canon. Aucun nouveau correctif n'est
introduit par la rédaction de ce document.

\section{Sony : ARW, variantes de compression et prudence spectrale}
\subsection{Ce qui est déjà documentable}
Les fichiers ARW relèvent d'une famille de structures
et de codecs Sony, avec métadonnées et aperçus.
Le code de RawSpeed fournit des exemples de décodage
de certaines variantes ARW, mais le nom d'une classe
ne démontre pas le support de tous les modes Sony~\cite{rawspeed,sonyarw2}.
Un lecteur doit reconnaître le modèle et la variante
avant d'interpréter profondeur, emballage et compression.

\subsection{Modes du Sony Alpha 7 IV : un exemple précis}
Le manuel Sony du 7 IV décrit des choix non compressé,
compressé sans perte L/M/S et compressé. Il précise que
M et S réduisent la résolution par rapport à L, et
mentionne une condition de firmware pour ces choix~\cite{sonymodes}.
Il ne faut pas lire « compression sans perte M » comme
« toutes les mesures de pleine résolution sont conservées » :
la compression peut être sans perte sur une représentation
déjà réduite spatialement.

Cette distinction s'applique bien au-delà de Sony :
la garantie d'un codec porte sur son entrée, pas sur
toutes les opérations antérieures de la chaîne caméra.
Pour un protocole rigoureux, choisir la pleine résolution
et documenter le mode, sans supposer qu'un autre modèle
dispose des mêmes menus.

\subsection{Un exemple ARW2 de codage par petits groupes}
Dans la variante illustrée par le décodeur SonyArw2
de RawSpeed, un groupe de seize échantillons consomme
128 bits, avec extrema, positions et valeurs résiduelles
compactes. Le code impose des invariants et reconstruit
les valeurs selon la plage du groupe~\cite{sonyarw2}.
Ce mécanisme est un exemple concret de codage adaptatif
à la dynamique locale, pas une description de tous les ARW.

Il explique pourquoi des erreurs ou compromis peuvent
être corrélés à des groupes et à de forts contrastes.
Les essais doivent inclure arêtes lumineuses, ombres
remontées et textures fines, et pas seulement des
aplats correctement exposés. Aucune estimation de
qualité Sony globale n'est déduite ici d'un seul codec.

\subsection{Couleur et futurs tests Sony dans IPR}
Une matrice Canon, ses gains et son profil ne peuvent
pas être réutilisés parce qu'un Sony aurait lui aussi
un CFA Bayer. Les réponses spectrales, calibrations,
normalisations et rendus diffèrent. Le futur corpus
Sony doit contenir un RAW plein format non réduit,
son mode exact, un rendu neutre de référence clairement
nommé et, idéalement, une charte sous éclairages connus.

\paragraph{À compléter avec des fichiers réels.} Inventorier
modèle/firmware, modes compressés et sans perte, géométrie
native, profondeur utile, niveaux, orientation, gains,
profil et effet des styles Sony. Cette section donne
une base documentée ; elle ne revendique aucune validation
de fichiers Sony dans la campagne Canon actuelle.

\section{Nikon : NEF et coexistence de procédés}
\subsection{NEF ne décrit pas un codec unique}
Les NEF utilisent des organisations de la famille TIFF
avec données Nikon et variantes de compression. Le
décodeur Nikon de RawSpeed illustre notamment des
tables, prédictions et décodages de différences pour
les variantes qu'il traite~\cite{nikondecoder}.
Ces mécanismes ne doivent pas être appliqués à tout
fichier Nikon sur le seul nom de marque.

\subsection{Exemple Z6/Z7 : profondeur et réversibilité}
Le manuel Nikon Z6/Z7 distingue compression sans
perte, compression non réversible et non compressé,
ainsi que des choix 12 ou 14 bits~\cite{nikonmodes}.
Les deux axes sont distincts : un RAW 14 bits avec
pertes et un RAW 12 bits sans perte ne se comparent
pas avec une seule étiquette « qualité ».

Le nombre de bits indique une capacité de codage,
pas à lui seul la dynamique effective du capteur.
Un choix de mode peut aussi modifier vitesse et
compatibilité ; il faut consulter le manuel du
modèle et son firmware au lieu d'extrapoler celui
d'un Z6 ou d'un Z7 à un Z9.

\subsection{Exemple Z6 III : haute efficacité}
Le manuel du Z6 III distingue compression sans
perte et deux variantes haute efficacité, avec
des tailles et objectifs de qualité différents~\cite{nikonefficiency}.
Le nom marketing « haute efficacité » ne suffit
pas à identifier une structure binaire ni à
garantir la réversibilité. La prise en charge
doit être vérifiée pour la bibliothèque et sa
version réellement utilisées.

Une bibliothèque peut ouvrir les métadonnées
et l'aperçu sans savoir développer un mode
de compression nouveau. Le test pertinent
demande les pixels natifs et leur cohérence,
pas seulement une icône visible dans l'explorateur.

\subsection{Picture Control et futurs tests Nikon}
Les réglages de rendu Nikon, comme les Picture
Controls, doivent être distingués de la mosaïque,
des métadonnées et des profils tiers. Pour une
comparaison, préciser si la référence est un
JPEG boîtier, un développement NX Studio ou
un développement Lightroom avec profil défini.

\paragraph{À compléter avec des fichiers réels.} Tester
au moins les modes pertinents du modèle retenu,
la profondeur, les dimensions, les aperçus,
l'orientation, les gains et la calibration.
Les formats de certaines gammes, comme NRW,
doivent faire l'objet d'une vérification séparée ;
ils ne sont pas couverts par une recette NEF
automatique supposée universelle.

\section{Synthèse inter-marques et limites de connaissance}
\begin{longtable}{p{3.2cm}p{4cm}p{7.5cm}}
\toprule Question & Invariant utile & Variable à vérifier \\ \midrule
Qu'est-ce qui est mesuré ? & Signal du capteur & CFA, profondeur, prétraitements et lecture. \\
Comment le fichier s'organise ? & Index et charges utiles & TIFF, boîtes, versions et offsets. \\
Que fait le codec ? & Reconstruit une représentation & Prédiction, transformée, pertes et mode réduit. \\
Comment retrouver les couleurs ? & Calibration nécessaire & Matrices, gains, tables et illuminants. \\
Quel rendu veut-on ? & Choix explicite & Style boîtier, profil tiers ou cible physique. \\
Peut-on faire confiance au décodage ? & Tests reproductibles & Modèle, firmware, mode et version bibliothèque. \\
\bottomrule
\end{longtable}
Le bon investissement logiciel consiste à généraliser
les contrats et les contrôles, pas à généraliser des
constantes issues d'un seul boîtier. Les informations
manquantes doivent rester identifiées comme telles.

\part{Architecture logicielle, diagnostic et travaux pratiques}
\section{Chaîne d'encodage et chaîne de décodage}
\subsection{Encoder une image développée}
Un export doit définir son contrat avant de choisir un
codec : dimensions, orientation, espace couleur, fonction
de transfert, profondeur, alpha et métadonnées. Une chaîne
conceptuelle correcte est : terminer les traitements dans
leur domaine approprié, effectuer le rendu de sortie,
convertir les couleurs, appliquer l'encodage de transfert,
quantifier, compresser puis écrire les structures et les
métadonnées cohérentes.

Le codec peut effectuer lui-même certaines de ces opérations.
Il faut connaître son API pour ne pas les appliquer deux
fois. Une conversion RGB vers YCbCr peut être gérée par
l'encodeur JPEG ; une bibliothèque TIFF peut attendre
des octets déjà organisés comme les tags le décrivent.

L'écriture devrait utiliser un fichier temporaire dans
un emplacement approprié, vérifier les erreurs puis
publier le résultat complet. Une taille non nulle ne
prouve pas que l'écriture est terminée. En cas de conflit
de nom, une stratégie d'incrémentation doit tester le
nom final complet et préserver les fichiers existants.
L'atomicité exacte dépend du système de fichiers et de
l'opération de publication choisie.

\subsection{Décoder une image rendue}
Le lecteur vérifie la structure, choisit le codec,
reconstruit les composantes, interprète leur type,
leur couleur et leur alpha, applique l'orientation
si demandée et produit un tampon interne. La conversion
vers l'écran vient ensuite. Il faut savoir si la
bibliothèque réalise déjà orientation, gestion ICC,
conversion de profondeur ou prémultiplication.

Un contrat comme « RGBA float32 linéaire sRGB, alpha droit,
orientation appliquée » est beaucoup plus utile que
« une image OpenCV ». Les erreurs inter-modules viennent
souvent d'informations restées implicites.

\section{Mathématiques de traitement : une carte pour se repérer}
\subsection{Convolution et séparation fréquentielle}
Un filtre discret linéaire peut s'écrire
\begin{equation}
J[x,y]=\sum_{u,v}K[u,v]I[x-u,y-v].
\end{equation}
Un noyau de moyenne réduit des variations rapides et
du bruit, mais aussi des détails. Un noyau passe-haut
réagit aux contours et au bruit. Les conditions aux
bords, répétition, symétrie ou remplissage, font partie
de l'algorithme et peuvent créer des différences de test.

La transformée de Fourier discrète 2D est
\begin{equation}
F[k,\ell]=\sum_{x=0}^{W-1}\sum_{y=0}^{H-1}
I[x,y]e^{-2\pi i(kx/W+\ell y/H)}.
\end{equation}
Elle décompose des motifs en fréquences globales ;
les ondelettes ajoutent une localisation et des
échelles. FFT désigne un algorithme rapide pour
calculer une transformée, pas un filtre en soi.

\subsection{Accentuation et débruitage}
Une accentuation classique utilise
$J=I+a(I-G_\sigma*I)$ : elle ajoute une différence
entre image et version floutée. Le paramètre $a$
augmente aussi le bruit et peut créer des halos.
L'ordre débruitage/netteté change donc le résultat.

Le débruitage est un problème d'estimation : retrouver
une image plausible à partir de mesures bruitées.
Un objectif simplifié est
\begin{equation}
\widehat I=\arg\min_J\left[
\frac{1}{2\sigma^2}\|Y-J\|_2^2+\lambda\mathcal R(J)\right].
\end{equation}
Le premier terme attache la solution aux données,
le second impose une régularité ou un a priori.
Les méthodes non locales, bayésiennes ou apprises
diffèrent notamment par cet a priori et leur modèle
de bruit. Un bruit de photons dépend du signal ;
un modèle gaussien constant n'est pas toujours suffisant.
Les cours 05, 08 et 10 développent les outils liés.

\subsection{Géométrie et rééchantillonnage}
Une transformation affine s'écrit en coordonnées
homogènes :
\begin{equation}
\begin{pmatrix}x'\\y'\\1\end{pmatrix}=
\begin{pmatrix}a&b&t_x\\c&d&t_y\\0&0&1\end{pmatrix}
\begin{pmatrix}x\\y\\1\end{pmatrix}.
\end{equation}
Une homographie remplace aussi la dernière ligne,
puis nécessite une division par la troisième
coordonnée. Pour produire une grille de sortie
sans trous, on utilise souvent la transformation
inverse : pour chaque pixel cible, chercher sa
position source et interpoler.

Un recadrage entier peut simplement changer les
bornes. Une rotation arbitraire, une correction
de perspective ou de distorsion impose généralement
un rééchantillonnage. Les conventions de centre
de pixel et le filtre doivent être identiques
pour comparer deux implémentations au pixel près.

\section{Aperçus, proxies, caches et performances}
\subsection{Pourquoi une image peut apparaître en deux temps}
Un aperçu intégré est rapide à afficher, mais représente
un rendu déjà calculé, parfois dans une résolution
et un profil différents. Le moteur peut ensuite
remplacer cet aperçu par son développement. Ce
changement ne signifie pas nécessairement qu'un
filtre utilisateur a été appliqué.

Un proxy est une représentation réduite choisie
pour l'interaction. Il peut accélérer le zoom
global et les réglages, mais ne suffit pas à
juger du bruit ou de la netteté à l'échelle native.
Un zoom 100\,\% du proxy n'est pas le 100\,\%
du RAW pleine résolution.

\subsection{Ce qu'une clé de cache doit distinguer}
Une clé robuste inclut l'identité de la source,
le niveau de résolution, l'orientation, les
paramètres de décodage, la calibration ou son
identité, la recette de rendu et la version
de l'étape concernée. Sinon, changer de profil
peut réutiliser une image développée avec l'ancien.

Un cache borné LRU élimine les entrées les moins
récemment utilisées selon un budget explicite.
Le budget doit être en octets ou en coût estimé,
pas seulement en nombre d'images de tailles
potentiellement différentes. Précharger des
vignettes ne justifie pas de garder indéfiniment
tous les RAW développés.

\subsection{Mesurer l'ouverture}
Chronométrer séparément : recherche des profils,
lecture, décompression, dématriçage, transformations,
copie mémoire, transfert GPU et affichage. Comparer
le premier passage, avec caches froids, aux suivants,
avec caches chauds. Une mesure globale de 30 secondes
ne permet pas de désigner le dématriçage ou le disque
comme responsable sans autre observation.

Les en-têtes et métadonnées peuvent être lus sans
développement complet. Une vignette peut provenir
du fichier ou d'un cache. L'architecture doit
prévoir annulation et priorité des demandes :
une image quittant la vue ne doit pas monopoliser
le calcul au détriment de celle que l'on ouvre.

\section{Sécurité, robustesse et archivage}
\subsection{Un fichier image est une entrée non fiable}
Valider tailles, offsets, profondeur, nombre de
canaux, ratio de décompression et profondeur
d'imbrication. Limiter mémoire et durée de calcul.
Une bombe de décompression est un fichier compact
qui demande énormément de mémoire après expansion.
Un graphe de répertoires peut contenir une boucle.
Un profil embarqué peut lui aussi être malformé.

Les CRC n'empêchent pas un attaquant de fabriquer
un fichier cohérent et malveillant. Les bibliothèques
doivent être maintenues, testées sur corpus variés
et éventuellement isolées selon le contexte. Le
fuzzing explore des entrées modifiées pour trouver
des plantages ou comportements indéfinis ; il
complète les tests de conformité, sans les remplacer.

\subsection{Vie privée et provenance}
Un export peut révéler lieu, date, numéro de série,
nom de propriétaire et parfois davantage de contexte
qu'attendu. Une option de suppression des métadonnées
doit distinguer données privées et informations
nécessaires à l'affichage, telles qu'un profil ICC.
Supprimer aveuglément tous les tags peut modifier
l'interprétation de l'image.

\subsection{Archivage durable}
Conserver original, sauvegardes vérifiées, hashes,
recette, profils autorisés et versions significatives.
Un rendu TIFF peut préserver une apparence stabilisée,
mais ne remplace pas toutes les possibilités du RAW.
Un DNG peut améliorer l'interopérabilité sans garantir
que chaque métadonnée constructeur aura été conservée
ou interprétée de manière identique. Tester une
restauration est plus probant que constater qu'une
copie existe.

\section{Diagnostic : quels nombres comparer ?}
\subsection{Trois niveaux de test}
\begin{enumerate}
\item \textbf{Binaire :} signature, structure, offsets,
CRC éventuels, hash et longueur. Cela vérifie le fichier.
\item \textbf{Échantillons :} dimensions, type,
orientation, tableau décodé et round-trip. Cela vérifie
les données reconstruites.
\item \textbf{Rendu :} profil, fonction de transfert,
géométrie, luminance, couleurs et régions. Cela vérifie
l'interprétation et les traitements.
\end{enumerate}
Un round-trip sans perte doit conserver exactement
les échantillons du domaine déclaré. Un rendu RAW
avec dématriçage n'est pas un round-trip réversible
vers la mosaïque d'origine.

\subsection{MAE, RMSE et PSNR}
Pour $N$ échantillons comparables $a_i,b_i$,
\begin{align}
\mathrm{MAE}&=\frac1N\sum_i|a_i-b_i|,\\
\mathrm{RMSE}&=\sqrt{\frac1N\sum_i(a_i-b_i)^2},\\
\mathrm{PSNR}&=20\log_{10}\frac{L_{\max}}{\mathrm{RMSE}}.
\end{align}
La plage $L_{\max}$ et le domaine, codes sRGB ou
valeurs linéaires par exemple, doivent être précisés.
Une erreur nulle donne un PSNR infini dans cette
définition ; il faut le représenter explicitement
dans les rapports plutôt que provoquer une division
par zéro.

Un décalage d'un pixel sur un contour peut produire
une forte MAE sans grande différence de couleur
locale. Une moyenne globale peut masquer une
dégradation importante sur la peau ou les ombres.
Mesurer des régions fixes et visualiser les erreurs.

\subsection{Luminance et erreur colorimétrique}
La luminance relative $Y$ se calcule à partir des
composantes linéaires avec les poids de l'espace,
pas directement avec les codes non linéaires.
Pour des moyennes strictement positives,
\begin{equation}
\Delta e_Y=\log_2\frac{\overline Y_{\text{référence}}}
{\overline Y_{\text{test}}}
\end{equation}
décrit un rapport en EV. Cela ne prouve pas qu'un
curseur d'exposition suffira à corriger le rendu.

La distance Euclidienne CIELAB, notée $\Delta E_{76}$,
n'est pas uniformément perceptuelle. CIEDE2000
pondère et couple différemment les écarts ; sa
mise en œuvre doit être testée avec des paires
de référence, comme celles de Sharma et collègues~\cite{ciedeimages}.
Un $\Delta E$ exige une conversion colorimétrique
cohérente et ne constitue pas seul une évaluation
du détail ou du bruit.

\subsection{Clipping et fractions sombres}
« Au moins un canal proche de zéro » n'est pas
« pixel noir ». Une couleur rouge pure possède
deux canaux nuls sans être noire. Il faut distinguer
fraction de canaux, fraction de pixels avec un
canal écrêté, fraction tous canaux écrêtés et
seuil de luminance. Pour une sortie 16 bits,
indiquer explicitement un seuil comme $1/65535$.

Le bruit estimé par MAD est la médiane des écarts
absolus à une médiane. Sur un coefficient gaussien
centré, $\sigma\approx\mathrm{MAD}/0{,}67448975$.
Il faut choisir une zone ou une transformation
qui sépare suffisamment le bruit de la texture.
Une faible dispersion après flou peut seulement
signifier que l'on a détruit du détail.

\section{Travaux pratiques reproductibles}
\subsection{Lancer le laboratoire autonome}
Depuis la racine du dépôt, avec Python 3 :
\begin{verbatim}
python docs/scholar/references/images-exemples.py
\end{verbatim}
Sans option, il affiche des calculs et lance ses
autotests, sans créer de fichier. Il n'utilise
que la bibliothèque standard Python. Pour fabriquer
trois images synthétiques dans un nouveau dossier :
\begin{verbatim}
python docs/scholar/references/images-exemples.py \
  --output-dir build-scholar-images-tp
\end{verbatim}
La continuation par antislash ci-dessus est celle
d'un shell POSIX ; sous PowerShell, écrire la
commande sur une ligne ou utiliser un accent grave.
Le dossier existant est refusé. Les photographies
du corpus et les profils Adobe ne sont pas modifiés.

Les trois fichiers montrent les mêmes deux pixels
rouge et vert dans PPM, PNG et TIFF. Le programme
compare les pixels de son PNG avec ceux d'origine,
teste les cinq filtres et vérifie les offsets TIFF.
Son décodeur PNG est volontairement limité à RGB
8 bits non entrelacé et refuse les autres variantes.
Il ne faut pas le brancher sur la production.

\subsection{Observer les octets}
\begin{verbatim}
python docs/scholar/references/images-exemples.py \
  --inspect build-scholar-images-tp/deux-pixels.tiff
\end{verbatim}
Cette option lit seulement les 64 premiers octets
et la taille. Comparer la signature au tableau
TIFF de la partie IV. Ouvrir ensuite les trois
fichiers dans un lecteur compatible et vérifier
qu'ils montrent bien deux pixels : rouge puis vert.
Une application peut agrandir ces pixels par
interpolation ; pour observer leurs codes, préférer
une lecture numérique ou un zoom sans lissage.

\subsection{Lire les métadonnées d'une photographie réelle}
Avec ExifTool installé, en lecture seule :
\begin{verbatim}
exiftool -G1 -a -s photo.CR3
exiftool -G1 -a -s photo.tiff
\end{verbatim}
Adapter l'extension à Sony ou Nikon. Relever les
dimensions, modèle, ISO, orientation, profondeur,
profil, compression, gains et informations de
développement disponibles. Les doublons sont
volontairement affichés avec leurs groupes.
Le cours ne recommande aucune commande d'écriture
de métadonnées sur les originaux.

\subsection{Éprouver la réversibilité}
Le laboratoire compacte $[\texttt{ABC},\texttt{123}]$
sur 12 bits et obtient \texttt{AB C1 23}. Deux
échantillons font 24 bits, donc trois octets.
Modifier un seul bit et observer quelle mesure
change. Ce TP illustre un empaquetage MSB simple,
pas le format natif d'un CR3 ou d'un ARW.

Pour la compression, comparer l'encodage d'une
plage uniforme et de données pseudo-aléatoires.
Comparer les tailles de fichiers, puis vérifier
les échantillons décodés. Ne pas prendre le ratio
de compression comme un score de qualité visuelle.

\subsection{Observer une perte sans logiciel photographique}
Calculer la DCT d'un bloc constant 12 : le coefficient
DC vaut 96. Quantifier ce coefficient avec $Q=20$
donne 5, puis une reconstruction DC=100. Le bloc
inverse devient 12,5 partout, avant les arrondis
de sortie. L'erreur de 0,5 n'est pas un défaut
de la transformée inverse : elle a été introduite
par la quantification.

\subsection{Protocole comparatif RAW minimal}
Conserver un original RAW, un aperçu extrait si
disponible et deux développements explicitement
paramétrés. Désactiver ou harmoniser l'optique,
le débruitage, la netteté et les réglages ; choisir
le même profil cible et documenter As Shot.
Exporter en TIFF sRGB 16 bits sans changement
de dimensions. Vérifier métadonnées et géométrie
avant de calculer MAE et différences Lab.
Ce protocole rapproche les conditions ; il ne
garantit pas que les moteurs ont le même rendu
de base ni qu'ils sont physiquement fidèles.

\section{Exercices corrigés}
\subsection*{1. Estimer la mémoire}
\emph{Question.} Une photo $8000\times6000$ est chargée
en RGB uint16 puis en RGBA float32. Quelles tailles ?
\emph{Correction.} Il y a 48 millions de pixels.
RGB uint16 : $48\times10^6\times6=288$ MB.
RGBA float32 : $48\times10^6\times16=768$ MB.
Garder les deux représente 1056 MB avant caches
et intermédiaires. La taille compressée n'intervient
pas dans ce calcul.

\subsection*{2. Lire un entier}
\emph{Question.} Que vaut \texttt{00 10} sur 16 bits ?
\emph{Correction.} Big-endian : 16. Little-endian :
4096. Le fichier doit annoncer sa convention.

\subsection*{3. Calculer une adresse}
\emph{Question.} RGB uint16, stride 40, pixel $(3,2)$,
canal 1, origine $A_0$ ?
\emph{Correction.} $A_0+2\times40+3\times3\times2+
1\times2=A_0+100$. Ne pas remplacer le stride
par une largeur supposée.

\subsection*{4. Normaliser un RAW}
\emph{Question.} $D=4096$, $B=1024$, $W=15360$.
\emph{Correction.} $(4096-1024)/(15360-1024)
=3072/14336=3/14\approx0{,}214286$.
Utiliser 65535 au dénominateur parce que le
tampon est uint16 serait une erreur.

\subsection*{5. Appliquer un EV}
\emph{Question.} $L=0{,}18$, réglage $+2$ EV.
\emph{Correction.} $L'=0{,}72$. Encoder seulement
ensuite pour une sortie sRGB ; multiplier
directement le code 0,46136 par quatre écrêterait
abusivement le résultat.

\subsection*{6. Expliquer un hash différent}
\emph{Question.} Deux PNG donnent les mêmes pixels,
mais ont des hashes différents. Est-ce une perte ?
\emph{Correction.} Pas nécessairement. Filtres,
compression, chunks ou métadonnées peuvent différer.
Il faut distinguer identité du fichier et égalité
de ses pixels interprétés.

\subsection*{7. Interpréter le clipping}
\emph{Question.} Le pixel linéaire $(1,0,0)$ est-il noir ?
\emph{Correction.} Non. Il est saturé en rouge dans
la représentation considérée, avec deux canaux
nuls. Un compteur « au moins un canal bas »
ne mesure donc pas le nombre de pixels noirs.

\subsection*{8. Comparer deux TIFF}
\emph{Question.} Même scène, même taille, forte MAE.
Que vérifier avant de corriger les couleurs ?
\emph{Correction.} Orientation, recadrage d'un
pixel, optique, dématriçage, profil ICC, fonction
de transfert, profondeur et normalisation, puis
réglages et rendu de base. Une taille identique
ne démontre pas l'alignement.

\subsection*{9. Comprendre « sans perte »}
\emph{Question.} Un RAW réduit est compressé sans
perte. Peut-on récupérer la résolution originale ?
\emph{Correction.} Pas en général. Le codec peut
conserver exactement une représentation déjà
réduite. La garantie porte sur son entrée.

\subsection*{10. Identifier une fausse conversion}
\emph{Question.} Renommer TIFF en DNG ou assigner
sRGB à des codes Adobe RGB convertit-il les données ?
\emph{Correction.} Non. Le premier change un nom,
le second une interprétation. Une conversion
nécessite des transformations et structures
cohérentes avec le résultat.

\subsection*{11. Diagnostiquer une rotation double}
\emph{Question.} Un export portrait apparaît
tourné encore de 90 degrés. Cause possible ?
\emph{Correction.} Les pixels ont pu être
réorientés sans remettre l'orientation à l'identité.
Vérifier le tableau et la métadonnée séparément.

\subsection*{12. Distinguer précision et dynamique}
\emph{Question.} Convertir un JPEG 8 bits en
TIFF 16 bits récupère-t-il les hautes lumières
perdues ?
\emph{Correction.} Non. Cela augmente la précision
disponible pour les traitements futurs, sans
reconstituer les valeurs déjà écrêtées ou quantifiées.

\section{Glossaire}
\begin{longtable}{p{3.4cm}p{11.2cm}}
\toprule Terme & Définition dans ce cours \\ \midrule
ADC / CAN & Convertisseur analogique-numérique. \\
Aliasing & Repliement de fréquences lié à un échantillonnage insuffisant. \\
Alpha & Information de couverture/opacité ; droit ou prémultiplié. \\
Aperçu & Image déjà développée, souvent embarquée et réduite. \\
Bit / octet & Chiffre binaire / groupe de huit bits. \\
Bit depth & Nombre de bits de représentation par échantillon. \\
Black level & Offset associé au signal noir dans une convention de codes. \\
BMFF & Famille d'organisation de fichiers média en boîtes. \\
Canal & Champ de valeurs : couleur, alpha, profondeur ou autre grandeur. \\
CFA & Mosaïque de filtres colorés du capteur. \\
Chroma & Composantes décrivant une différence colorée dans un modèle donné. \\
Clipping & Écrêtage de valeurs à une limite. \\
Codec & Règles et implémentation de codage/décodage d'une représentation. \\
Conteneur & Structure organisant flux, propriétés et relations. \\
CRC & Contrôle d'intégrité destiné notamment aux erreurs accidentelles. \\
DCP & Profil de caméra de la famille DNG ; matrices et rendu possibles. \\
DCT & Transformée en cosinus discrète. \\
Dématriçage & Estimation de composantes couleur absentes du CFA. \\
DEFLATE & Compression combinant références LZ77 et codage Huffman. \\
Dithering & Signal ajouté pour rendre certaines erreurs de quantification moins structurées. \\
DNG & Format de négatif numérique fondé sur TIFF et enrichi de calibrations. \\
Endianness & Ordre des octets d'une valeur multioctet. \\
Entropie & Quantité moyenne d'information dans un modèle probabiliste. \\
EOTF / OETF & Relations de transfert vers la lumière d'affichage / vers un signal encodé. \\
EV & Unité logarithmique d'exposition ; un EV correspond à un facteur deux. \\
Exif & Conventions de données photographiques et structures associées. \\
FFT & Famille d'algorithmes rapides de transformée de Fourier discrète. \\
Gamut & Ensemble des couleurs représentables ou reproductibles. \\
Gain & Facteur multiplicatif, à distinguer d'une nouvelle acquisition de photons. \\
HDR / SDR & Contextes de dynamique étendue / standard, à préciser selon l'étape. \\
ICC & Famille de profils et contrats de gestion colorimétrique. \\
IFD & Répertoire de champs d'image dans TIFF et familles apparentées. \\
Illuminant & Description spectrale ou normalisée d'un éclairage de référence. \\
IPTC & Conventions éditoriales, descriptives et administratives des médias. \\
JPEG & Famille de procédés de codage ; le mode courant utilise DCT et pertes. \\
LCP & Profil de corrections d'objectif, distinct d'un profil couleur caméra. \\
Linéaire & Proportionnel à une grandeur explicitement définie. \\
Luma / luminance & Signal pondéré souvent non linéaire / grandeur colorimétrique Y. \\
LUT & Table de correspondance, éventuellement multidimensionnelle. \\
MAD & Médiane des écarts absolus à une médiane. \\
MakerNote & Bloc de métadonnées propre au constructeur. \\
Métamérisme & Réponses identiques malgré des spectres différents. \\
MTF / PSF & Transmission des contrastes / étalement spatial d'un point. \\
Offset & Position binaire relative à une origine ; ou décalage additif selon le contexte. \\
Optical-black & Zone de capteur masquée utilisée pour caractériser le noir. \\
PCS & Espace de connexion des profils d'une architecture ICC. \\
Photosite & Élément spatial de mesure du capteur. \\
Pixel & Position/élément d'une représentation matricielle, dont le sens dépend du stade. \\
Planar / interleaved & Composantes en plans séparés / entrelacées. \\
Profil & Description de transformation ; son rôle dépend du type ICC, DCP, etc. \\
Proxy & Représentation réduite pour l'interaction. \\
Quantification & Passage à un ensemble discret de valeurs. \\
RAW & Famille de données photographiques proches des mesures, pas un codec unique. \\
Résidu & Différence entre valeur et prédiction. \\
ROI & Région d'intérêt définie dans un repère et un domaine précis. \\
Round-trip & Encodage puis décodage avec comparaison au domaine initial. \\
Sidecar & Fichier associé portant paramètres ou métadonnées. \\
SNR & Rapport signal/bruit, avec conventions à préciser. \\
sRGB & Espace RGB normalisé avec primaires, blanc et transfert définis. \\
Stride / pitch & Distance mémoire entre deux débuts de ligne. \\
Tag & Champ identifié dans une structure de données. \\
Tone mapping & Adaptation du rendu tonal vers une capacité de sortie. \\
Tuile / bande & Bloc spatial rectangulaire / portion horizontale de stockage. \\
White level & Niveau de référence de saturation ou normalisation, selon le contexte. \\
XMP & Modèle extensible de métadonnées, souvent sérialisé en RDF/XML. \\
XYZ / Lab & Espaces colorimétriques, non simples synonymes de RGB. \\
\bottomrule
\end{longtable}

\section{Bibliographie : quoi lire, dans quel ordre ?}
Pour la structure binaire, commencer par les exemples de
ce cours puis les spécifications PNG, TIFF et Netpbm.
Pour la compression, suivre les calculs de prédiction et
DCT, puis T.81 et les RFC zlib/DEFLATE. Pour la couleur,
lire la fiche sRGB et l'architecture ICC avant les tags
avancés DNG. Pour les capteurs, EMVA fournit un cadre de
mesure plutôt qu'un simple classement commercial.

Les manuels Sony et Nikon de cette bibliographie sont
volontairement attachés à des modèles. Les ressources
Clévy et les décodeurs LibRaw/RawSpeed donnent accès aux
détails des RAW, mais il faut lire leurs limites et leurs
licences. Les branches de code en ligne évoluent : conserver
un commit et les fichiers de test pour une investigation
reproductible de production.

La ressource historique CR2 \url{http://lclevy.free.fr/cr2/}
reste intéressante, mais son accès direct a échoué pendant
la recherche. Le dépôt de l'auteur \texttt{lclevy/libcraw2}
est donc proposé comme point d'entrée alternatif, sans
prétendre avoir relu une page inaccessible. Certains
catalogues de tags ExifTool limitent également l'accès
automatisé ; leur index et les outils locaux permettent
de poursuivre la lecture.

\resultat{Pour comprendre une image, toujours demander : que
mesurent les nombres, comment sont-ils rangés, quelles transformations
ont déjà été faites, et quelles informations faut-il encore appliquer
avant de les afficher ?}


\clearpage
\bibliographystyle{plain}
\bibliography{references-images}
\end{document}
